
\Data\ID{1003020201}
\Updated{2021-04-26 03:04:57}
\Level{A}
\SubArea{Linear equations and inequalities}
\LANGen{
\Question{
Find the solution set of the following equation.
\[ 2-\frac{2-x}3=\frac x2+\frac{8-x}6 \]}
{\( \mathbb{R} \)}{\( \emptyset \)}{\( \{0\} \)}{\( \{4\} \)}{}{}}
\LANGpl{
\Question{
Wyznacz zbiór rozwiązań równania.
\[ 2-\frac{2-x}3=\frac x2+\frac{8-x}6 \]}
{\( \mathbb{R} \)}{\( \emptyset \)}{\( \{0\} \)}{\( \{4\} \)}{}{}}
\LANGcs{
\Question{
Určete množinu řešení rovnice.
\[ 2-\frac{2-x}3=\frac x2+\frac{8-x}6 \]}
{\( \mathbb{R} \)}{\( \emptyset \)}{\( \{0\} \)}{\( \{4\} \)}{}{}}
\LANGsk{
\Question{
Určte množinu riešení rovnice.
\[ 2-\frac{2-x}3=\frac x2+\frac{8-x}6 \]}
{\( \mathbb{R} \)}{\( \emptyset \)}{\( \{0\} \)}{\( \{4\} \)}{}{}}
\LANGes{
\Question{
Calcula el conjunto de soluciones de la ecuación.
\[ 2-\frac{2-x}3=\frac x2+\frac{8-x}6\text{.} \]}
{\( \mathbb{R} \)}{\( \emptyset \)}{\( \{0\} \)}{\( \{4\} \)}{}{}}
\End

\Data\ID{1003020202}
\Updated{2021-04-26 03:04:17}
\Level{A}
\SubArea{Linear equations and inequalities}
\LANGen{
\Question{
Assuming \( x\in\mathbb{N} \), find the solution set of the equation.
\[ 3x-\left[1+2\cdot(3x-2)\right]=9\]}
{\( \emptyset \)}{\( \mathbb{N} \)}{\( \{-2\} \)}{\( \left\{\frac{14}9\right\} \)}{}{}}
\LANGpl{
\Question{
Zakładając, że \( x\in\mathbb{N} \), wyznacz zbiór rozwiązań równania
\[ 3x-\left[1+2\cdot(3x-2)\right]=9\text{ .} \]}
{\( \emptyset \)}{\( \mathbb{N} \)}{\( \{-2\} \)}{\( \left\{\frac{14}9\right\} \)}{}{}}
\LANGcs{
\Question{
Určete množinu řešení rovnice pro \( x\in\mathbb{N} \)
\[ 3x-\left[1+2\cdot(3x-2)\right]=9\text{ .} \]}
{\( \emptyset \)}{\( \mathbb{N} \)}{\( \{-2\} \)}{\( \left\{\frac{14}9\right\} \)}{}{}}
\LANGsk{
\Question{
Určte množinu riešení rovnice pre  \( x\in\mathbb{N} \)
\[ 3x-\left[1+2\cdot(3x-2)\right]=9\text{ .} \]}
{\( \emptyset \)}{\( \mathbb{N} \)}{\( \{-2\} \)}{\( \left\{\frac{14}9\right\} \)}{}{}}
\LANGes{
\Question{
Suponiendo que \( x\in\mathbb{N} \), calcula el conjunto de soluciones de la ecuación.
\[ 3x-\left[1+2\cdot(3x-2)\right]=9\text{ .} \]}
{\( \emptyset \)}{\( \mathbb{N} \)}{\( \{-2\} \)}{\( \left\{\frac{14}9\right\} \)}{}{}}
\End

\Data\ID{1003020203}
\Updated{2021-04-26 03:04:09}
\Level{A}
\SubArea{Linear equations and inequalities}
\LANGen{
\Question{
Find the solution set of the following equation.
\[ \frac{x+5}2-\frac{x+1}4=\frac{22+2x}8\]}
{\( \emptyset \)}{\( \mathbb{R} \)}{\( \{4\} \)}{\( \{-4\} \)}{}{}}
\LANGpl{
\Question{
Wyznacz zbiór rozwiązań równania.
\[ \frac{x+5}2-\frac{x+1}4=\frac{22+2x}8 \]}
{\( \emptyset \)}{\( \mathbb{R} \)}{\( \{4\} \)}{\( \{-4\} \)}{}{}}
\LANGcs{
\Question{
Určete množinu řešení rovnice.
\[ \frac{x+5}2-\frac{x+1}4=\frac{22+2x}8 \]}
{\( \emptyset \)}{\( \mathbb{R} \)}{\( \{4\} \)}{\( \{-4\} \)}{}{}}
\LANGsk{
\Question{
Určte množinu riešení rovnice.
\[ \frac{x+5}2-\frac{x+1}4=\frac{22+2x}8 \]}
{\( \emptyset \)}{\( \mathbb{R} \)}{\( \{4\} \)}{\( \{-4\} \)}{}{}}
\LANGes{
\Question{
Calcula el conjunto de soluciones de la ecuación.
\[ \frac{x+5}2-\frac{x+1}4=\frac{22+2x}8\text{.} \]}
{\( \emptyset \)}{\( \mathbb{R} \)}{\( \{4\} \)}{\( \{-4\} \)}{}{}}
\End

\Data\ID{1003020204}
\Updated{2021-04-26 03:04:38}
\Level{A}
\SubArea{Linear equations and inequalities}
\LANGen{
\Question{
Assuming \( x\in\mathbb{Z} \), find the solution set of the equation.
\[ 3-6x+3\cdot\left\{x-\left[2-(x+1)\right]\right\}=0\]}
{\( \mathbb{Z} \)}{\( \emptyset \)}{\( \{1\} \)}{\( \{0\} \)}{}{}}
\LANGpl{
\Question{
Zakładając, że \( x\in\mathbb{Z} \), wyznacz zbiór rozwiązań równania
\[ 3-6x+3\cdot\left\{x-\left[2-(x+1)\right]\right\}=0\text{ .} \]}
{\( \mathbb{Z} \)}{\( \emptyset \)}{\( \{1\} \)}{\( \{0\} \)}{}{}}
\LANGcs{
\Question{
Určete množinu řešení rovnice pro \( x\in\mathbb{Z} \)
\[ 3-6x+3\cdot\left\{x-\left[2-(x+1)\right]\right\}=0\text{ .} \]}
{\( \mathbb{Z} \)}{\( \emptyset \)}{\( \{1\} \)}{\( \{0\} \)}{}{}}
\LANGsk{
\Question{
Pre \( x\in\mathbb{Z} \) určte množinu riešení rovnice
\[ 3-6x+3\cdot\left\{x-\left[2-(x+1)\right]\right\}=0\text{ .} \]}
{\( \mathbb{Z} \)}{\( \emptyset \)}{\( \{1\} \)}{\( \{0\} \)}{}{}}
\LANGes{
\Question{
Suponiendo que \( x\in\mathbb{Z} \), calcula el conjunto de soluciones de la ecuación.
\[ 3-6x+3\cdot\left\{x-\left[2-(x+1)\right]\right\}=0\text{ .} \]}
{\( \mathbb{Z} \)}{\( \emptyset \)}{\( \{1\} \)}{\( \{0\} \)}{}{}}
\End

\Data\ID{1003037301}
\Updated{2020-07-15 03:07:12}
\Level{A}
\SubArea{Linear equations and inequalities}
\LANGen{
\Question{
Five times the unknown number is greater by nine than one-half of the number. Find the number.}
{\( 2 \)}{\( -2 \)}{\( 1 \)}{\( 10 \)}{}{}}
\LANGpl{
\Question{
Pięciokrotność pewnej liczby jest większa o dziewięć niż połowa  tej liczby. Jaka to liczba?}
{\( 2 \)}{\( -2 \)}{\( 1 \)}{\( 10 \)}{}{}}
\LANGcs{
\Question{
Pětkrát zvětšené neznámé číslo je o devět větší než jeho polovina. Určete neznámé číslo.}
{\( 2 \)}{\( -2 \)}{\( 1 \)}{\( 10 \)}{}{}}
\LANGsk{
\Question{
Päťkrát zväčšené neznáme číslo je o deväť väčšie ako jeho polovica. Určte neznáme číslo.}
{\( 2 \)}{\( -2 \)}{\( 1 \)}{\( 10 \)}{}{}}
\LANGes{
\Question{
El quíntuple de un número desconocido es mayor en nueve unidades que la mitad del número desconocido.¿Cuál es el valor del número desconocido?}
{\( 2 \)}{\( -2 \)}{\( 1 \)}{\( 10 \)}{}{}}
\End

\Data\ID{1003037302}
\Updated{2020-07-15 03:07:12}
\Level{A}
\SubArea{Linear equations and inequalities}
\LANGen{
\Question{
The sum of three consecutive odd numbers is twenty greater than the smallest of these numbers. Find the smallest one of the three unknown odd numbers.}
{\( 7 \)}{\( 13 \)}{\( 21 \)}{\( 1 \)}{}{}}
\LANGpl{
\Question{
Suma trzech kolejnych liczb nieparzystych jest dwadzieścia razy większa niż najmniejsza z tych liczb. Podaj wartość najmniejszej z tych trzech liczb}
{\( 7 \)}{\( 13 \)}{\( 21 \)}{\( 1 \)}{}{}}
\LANGcs{
\Question{
Součet tří po sobě jdoucích lichých přirozených čísel je o dvacet větší než nejmenší z daných tří čísel. Určete nejmenší z dané trojice lichých čísel.}
{\( 7 \)}{\( 13 \)}{\( 21 \)}{\( 1 \)}{}{}}
\LANGsk{
\Question{
Súčet troch po sebe idúcich nepárnych prirodzených čísel je o dvadsať väčší než najmenšie z daných troch čísel. Určte najmenšie číslo z danej trojice nepárnych čísel.}
{\( 7 \)}{\( 13 \)}{\( 21 \)}{\( 1 \)}{}{}}
\LANGes{
\Question{
La suma de tres números impares consecutivos es veinte unidades mayor que el menor de los tres números. ¿Cuál es el valor del menor de los tres números?}
{\( 7 \)}{\( 13 \)}{\( 21 \)}{\( 1 \)}{}{}}
\End

\Data\ID{1003037303}
\Updated{2020-07-15 03:07:12}
\Level{A}
\SubArea{Linear equations and inequalities}
\LANGen{
\Question{
Two more than eight times of the unknown number is the same as one-third of the difference of one and the number. Find the number.}
{\( -\frac15 \)}{\( -\frac7{23} \)}{\( -\frac5{27} \)}{\( -\frac7{21} \)}{}{}}
\LANGpl{
\Question{
Dwa razy więcej niż osiem razy  nieznana liczba jest równe jednej trzeciej różnicy jeden i tej liczby. Jaka to liczba?}
{\( -\frac15 \)}{\( -\frac7{23} \)}{\( -\frac5{27} \)}{\( -\frac7{21} \)}{}{}}
\LANGcs{
\Question{
Osminásobek neznámého čísla zvětšený o dvě je totéž jako třetina rozdílu čísla jedna a neznámého čísla. Určete neznámé číslo.}
{\( -\frac15 \)}{\( -\frac7{23} \)}{\( -\frac5{27} \)}{\( -\frac7{21} \)}{}{}}
\LANGsk{
\Question{
Osemnásobok neznámeho čísla zväčšený o dva je to isté ako tretina rozdielu čísla jeden a neznámeho čísla.  Určte neznáme číslo.}
{\( -\frac15 \)}{\( -\frac7{23} \)}{\( -\frac5{27} \)}{\( -\frac7{21} \)}{}{}}
\LANGes{
\Question{
El óctuple de un número desconocido aumentado en dos equivale a un tercio de la diferencia entre 1 y del número desconocido. ¿Cuál es el valor del número desconocido?}
{\( -\frac15 \)}{\( -\frac7{23} \)}{\( -\frac5{27} \)}{\( -\frac7{21} \)}{}{}}
\End

\Data\ID{1003037304}
\Updated{2020-07-15 03:07:12}
\Level{A}
\SubArea{Linear equations and inequalities}
\LANGen{
\Question{
The product of two consecutive natural numbers is the same as the sum of the square of the smaller one and eight. Find the smaller number.}
{\( 8 \)}{\( 4 \)}{\( 2 \)}{There is not such a natural number.}{}{}}
\LANGpl{
\Question{
Iloczyn dwóch kolejnych liczb naturalnych jest taki sam jak suma kwadratu najmniejszej z tych liczb i osiem. Podaj wartość mniejszej liczby.}
{\( 8 \)}{\( 4 \)}{\( 2 \)}{Nie ma takiej liczby naturalnej.}{}{}}
\LANGcs{
\Question{
Součin dvou po sobě jdoucích přirozených čísel je stejný jako součet druhé mocniny menšího z daných přirozených čísel a čísla osm. Určete menší z přirozených čísel.}
{\( 8 \)}{\( 4 \)}{\( 2 \)}{Takové přirozené číslo neexistuje.}{}{}}
\LANGsk{
\Question{
Súčin dvoch po sebe idúcich prirodzených čísel je rovnaký ako súčet druhej mocniny menšieho z daných prirodzených čísel a čísla osem. Určte menšie číslo.}
{\( 8 \)}{\( 4 \)}{\( 2 \)}{Neexistuje také prirodzené číslo.}{}{}}
\LANGes{
\Question{
El producto de dos números naturales consecutivos equivale a la suma del cuadrado del menor y ocho. ¿Cuál es el valor del número más pequeño?}
{\( 8 \)}{\( 4 \)}{\( 2 \)}{No hay ningún número natural así.}{}{}}
\End

\Data\ID{1003037305}
\Updated{2020-07-15 03:07:12}
\Level{A}
\SubArea{Linear equations and inequalities}
\LANGen{
\Question{
Suppose we increase the numerator of five-ninth by an integer and we decrease the denominator of five-ninth by the same integer. We get the fraction with numerator six times greater than the denominator. Find the integer.}
{\( 7 \)}{\( 3 \)}{There is no such an integer.}{\( 1 \)}{}{}}
\LANGpl{
\Question{
Przypuśćmy, że zwiększymy licznik pięciu-dziewiątych o liczbę całkowitą i zmniejszymy mianownik pięciu-dziewiątych o tę samą liczbę całkowitą. Otrzymujemy ułamek z licznikiem sześć razy większy od mianownika. Znajdź tę liczbę całkowitą.}
{\( 7 \)}{\( 3 \)}{Nie ma takiej liczby całkowitej.}{\( 1 \)}{}{}}
\LANGcs{
\Question{
Jestliže čitatele zlomku pět devítin zvětšíme a jmenovatele zmenšíme o stejné přirozené číslo, dostaneme zlomek, jehož čitatel je šestkrát větší než jmenovatel. Kterým celým číslem upravujeme původní zlomek?}
{\( 7 \)}{\( 3 \)}{Takové celé číslo neexistuje.}{\( 1 \)}{}{}}
\LANGsk{
\Question{
Ak čitateľa zlomku päť devätin zväčšíme a menovateľa zmenšíme o rovnaké celé číslo, dostaneme zlomok, ktorého čitateľ je šesťkrát väčší než menovateľ.   Nájdi celé číslo.}
{\( 7 \)}{\( 3 \)}{Neexistuje také celé číslo.}{\( 1 \)}{}{}}
\LANGes{
\Question{
Si aumentamos el numerador de la fracción cinco novenos en un número entero y a la vez disminuimos el denominador con el mismo número entero, obtendremos la fracción cuyo numerador es seis veces mayor que su denominador. Encuentra el número entero.}
{\( 7 \)}{\( 3 \)}{No hay ningún número natural así.}{\( 1 \)}{}{}}
\End

\Data\ID{1003037306}
\Updated{2020-07-15 03:07:12}
\Level{A}
\SubArea{Linear equations and inequalities}
\LANGen{
\Question{
The ratio of two more than five times \( x \) and five is equal to the sum of \( x \) and \( a \). Determine \( a \).}
{\( 0.4 \)}{It is not possible to determine the number  \( a \) uniquely without knowing the number \( x \).}{\( 2 \)}{\( 10 \)}{}{}}
\LANGpl{
\Question{
Stosunek dwóch więcej niż pięć razy \( x \) i pięć jest równy sumie  \( x \) i  \( a \). Określ wartość \( a \).}
{\( 0{,}4 \)}{Niemożliwe jest określenie wartości  \( a \) specyficznie nie znając wartości liczby  \( x \).}{\( 2 \)}{\( 10 \)}{}{}}
\LANGcs{
\Question{
Poměr pětinásobku čísla \( x \) zvětšeného o dvě a pěti je roven součtu čísel \( x \) a \( a \). Určete \( a \).}
{\( 0{,}4 \)}{Není možné jednoznačně určit číslo \( a \) bez znalosti čísla \( x \).}{\( 2 \)}{\( 10 \)}{}{}}
\LANGsk{
\Question{
Pomer päťnásobku čísla \( x \) zväčšeného o dve a päť sa rovná súčtu čísel \( x \) a \( a \). Určte \( a \).}
{\( 0{,}4 \)}{Nie je možné určiť číslo \( a \) bez poznania čísla \( x \).}{\( 2 \)}{\( 10 \)}{}{}}
\LANGes{
\Question{
La proporción entre el quíntuple del número \( x \) aumentado en dos y cinco equivale a la suma de los números \( x \) y \( a \). Determina \( a \).}
{\( 0.4 \)}{No es posible exactamente determinar el número  \( a \) sin conocer el número \( x \).}{\( 2 \)}{\( 10 \)}{}{}}
\End

\Data\ID{1003047002}
\Updated{2020-07-15 03:07:12}
\Level{A}
\SubArea{Linear equations and inequalities}
\LANGen{
\Question{
We are given the equation \( 1-\frac{5-x}2=\frac x4 \).  Decide which of the following equations is equivalent to the given equation, i.e. which of the following equations was obtained from the given equation by equivalent transformations.}
{\( -6+2x=x \)}{\( -6-2x=x \)}{\( -9+2x=x \)}{\( -6-x=x \)}{}{}}
\LANGpl{
\Question{
Dane jest równanie  \( 1-\frac{5-x}2=\frac x4 \).  Zdecyduj, które z  poniższych równań jest równoważne z podanym wcześniej równaniem, tzn. które z poniższych równań powstało w wyniku przekształceń podanego równania.}
{\( -6+2x=x \)}{\( -6-2x=x \)}{\( -9+2x=x \)}{\( -6-x=x \)}{}{}}
\LANGcs{
\Question{
Je dána rovnice \( 1-\frac{5-x}2=\frac x4 \).  Která z následujících rovnic je s ní ekvivalentní, tj. vznikla ekvivalentní úpravou zadané rovnice?}
{\( -6+2x=x \)}{\( -6-2x=x \)}{\( -9+2x=x \)}{\( -6-x=x \)}{}{}}
\LANGsk{
\Question{
Je daná rovnica \( 1-\frac{5-x}2=\frac x4 \).  Ktorá z nasledujúcich rovníc je s ňou ekvivalentná, t.j. vznikla ekvivalentnými úpravami zadanej rovnice?}
{\( -6+2x=x \)}{\( -6-2x=x \)}{\( -9+2x=x \)}{\( -6-x=x \)}{}{}}
\LANGes{
\Question{
Dada la ecuación \( 1-\frac{5-x}2=\frac x4 \).  Determina cuál de las siguientes ecuaciones es equivalente a la ecuación dada.}
{\( -6+2x=x \)}{\( -6-2x=x \)}{\( -9+2x=x \)}{\( -6-x=x \)}{}{}}
\End

\Data\ID{1003047003}
\Updated{2020-07-15 03:07:12}
\Level{A}
\SubArea{Linear equations and inequalities}
\LANGen{
\Question{
We are given the equation  \( \frac{8x}{x+2}+\frac{12}{x+2}=\frac{2x}{x+2} \).  Decide which of the following equations has the different set of roots than the given equation has, i.e. choose the equation which is not equivalent to the given equation.}
{\( 8x+12=2x \)}{\( \frac{4x}{x+2}+\frac6{x+2}=\frac x{x+2} \)}{\( \frac{6x}{x+2}=-\frac{12}{x+2} \)}{\( \frac x{x+2}=-\frac2{x+2} \)}{}{}}
\LANGpl{
\Question{
Dane jest równanie \( \frac{8x}{x+2}+\frac{12}{x+2}=\frac{2x}{x+2} \). Zdecyduj, które z poniższych równań ma inny zbiór pierwiastków niż podane wcześniej równanie, tzn. wybierz równanie, które nie jest odpowiednikiem podanego równania.}
{\( 8x+12=2x \)}{\( \frac{4x}{x+2}+\frac6{x+2}=\frac x{x+2} \)}{\( \frac{6x}{x+2}=-\frac{12}{x+2} \)}{\( \frac x{x+2}=-\frac2{x+2} \)}{}{}}
\LANGcs{
\Question{
Je dána rovnice  \( \frac{8x}{x+2}+\frac{12}{x+2}=\frac{2x}{x+2} \). Z následujících rovnic vyberte tu, která má jinou množinu kořenů než zadaná rovnice, tj. není s danou rovnicí ekvivalentní.}
{\( 8x+12=2x \)}{\( \frac{4x}{x+2}+\frac6{x+2}=\frac x{x+2} \)}{\( \frac{6x}{x+2}=-\frac{12}{x+2} \)}{\( \frac x{x+2}=-\frac2{x+2} \)}{}{}}
\LANGsk{
\Question{
Je daná rovnica \( \frac{8x}{x+2}+\frac{12}{x+2}=\frac{2x}{x+2} \).  Z nasledujúcich rovníc vyberte tú, ktorá má inú množinu koreňov, než zadaná rovnica, t.j. nie je s danou rovnicou ekvivalentná.}
{\( 8x+12=2x \)}{\( \frac{4x}{x+2}+\frac6{x+2}=\frac x{x+2} \)}{\( \frac{6x}{x+2}=-\frac{12}{x+2} \)}{\( \frac x{x+2}=-\frac2{x+2} \)}{}{}}
\LANGes{
\Question{
Dada la ecuación  \( \frac{8x}{x+2}+\frac{12}{x+2}=\frac{2x}{x+2} \).  Determina cuál de las siguientes ecuaciones tiene un conjunto de raíces (soluciones) diferentes que la ecuación dada, es decir, no es equivalente a ella.}
{\( 8x+12=2x \)}{\( \frac{4x}{x+2}+\frac6{x+2}=\frac x{x+2} \)}{\( \frac{6x}{x+2}=-\frac{12}{x+2} \)}{\( \frac x{x+2}=-\frac2{x+2} \)}{}{}}
\End

\Data\ID{1103049401}
\Updated{2020-07-15 03:07:12}
\Level{A}
\SubArea{Linear equations and inequalities}
\IMG{1103049401.pdf}
\LANGen{
\Question{
Choose the equation whose graphical solution is in the picture.}
{\( \frac12x-3=-\frac52x-5 \)}{\( \frac12x-3=-2x-5 \)}{\( 2x-3=-\frac52 x-5 \)}{\( 2x-3=-2x-5 \)}{}{}}
\LANGpl{
\Question{
Wybierz równanie, którego wykres przedstawiono poniżej.}
{\( \frac12x-3=-\frac52x-5 \)}{\( \frac12x-3=-2x-5 \)}{\( 2x-3=-\frac52 x-5 \)}{\( 2x-3=-2x-5 \)}{}{}}
\LANGcs{
\Question{
Vyberte rovnici, jejíž grafické řešení je na obrázku.}
{\( \frac12x-3=-\frac52x-5 \)}{\( \frac12x-3=-2x-5 \)}{\( 2x-3=-\frac52 x-5 \)}{\( 2x-3=-2x-5 \)}{}{}}
\LANGsk{
\Question{
Vyberte rovnicu, ktorej grafické riešenie je na obrázku.}
{\( \frac12x-3=-\frac52x-5 \)}{\( \frac12x-3=-2x-5 \)}{\( 2x-3=-\frac52 x-5 \)}{\( 2x-3=-2x-5 \)}{}{}}
\LANGes{
\Question{
Elige la ecuación cuya solución gráfica está representada en la imagen.}
{\( \frac12x-3=-\frac52x-5 \)}{\( \frac12x-3=-2x-5 \)}{\( 2x-3=-\frac52 x-5 \)}{\( 2x-3=-2x-5 \)}{}{}}
\End

\Data\ID{1103049402}
\Updated{2020-07-15 03:07:10}
\Level{A}
\SubArea{Linear equations and inequalities}
\LANGen{
\Question{
Choose the picture which represents the graphical solution of the following equation.
\[ \frac23 x+\frac13=1-\frac32 x \]}
{}{}{}{}{}{}}
\LANGpl{
\Question{
Wybierz rysunek przedstawiający graficzne rozwiązanie podanego równania.
\[ \frac23 x+\frac13=1-\frac32 x \]}
{}{}{}{}{}{}}
\LANGcs{
\Question{
Vyberte obrázek, který představuje grafické řešení rovnice 
\[ \frac23 x+\frac13=1-\frac32 x \]}
{}{}{}{}{}{}}
\LANGsk{
\Question{
Vyberte obrázok, ktorý predstavuje grafické riešenie rovnice.
\[ \frac23 x+\frac13=1-\frac32 x \]}
{}{}{}{}{}{}}
\LANGes{
\Question{
NA}
{}{}{}{}{}{}}

\IMGanswer{1103049402a.pdf}
\IMGanswer{1103049402b.pdf}
\IMGanswer{1103049402c.pdf}
\IMGanswer{1103049402d.pdf}\End

\Data\ID{1103049403}
\Updated{2020-07-15 03:07:10}
\Level{A}
\SubArea{Linear equations and inequalities}
\LANGen{
\Question{
Choose the picture which represents the graphical solution of the following equation. 
\[ 1-\frac{5+2x}3=3 \]}
{}{}{}{}{}{}}
\LANGpl{
\Question{
Wybierz rysunek przedstawiający graficzne rozwiązanie podanego równania. 
\[ 1-\frac{5+2x}3=3 \]}
{}{}{}{}{}{}}
\LANGcs{
\Question{
Vyberte obrázek, který představuje grafické řešení rovnice.
\[ 1-\frac{5+2x}3=3 \]}
{}{}{}{}{}{}}
\LANGsk{
\Question{
Vyberte obrázok, ktorý predstavuje grafické riešenie rovnice. 
\[ 1-\frac{5+2x}3=3 \]}
{}{}{}{}{}{}}
\LANGes{
\Question{
NA}
{}{}{}{}{}{}}

\IMGanswer{1103049403a.pdf}
\IMGanswer{1103049403b.pdf}
\IMGanswer{1103049403c.pdf}
\IMGanswer{1103049403d.pdf}\End

\Data\ID{1103049404}
\Updated{2020-07-15 03:07:10}
\Level{A}
\SubArea{Linear equations and inequalities}
\LANGen{
\Question{
In addition to solving equation \( ax+b=cx+d \) algebraically, you can also solve it graphically. When the lines \( y=ax+b \) and \( y=cx+d \) are graphed, you look for the intersection of these lines. In the pictures below are graphed lines \( y=ax+b \) and \( y=cx+d\).  Choose the picture in which the equation \( ax+b=cx+d \) has only one non-negative solution.}
{}{}{}{}{}{}}
\LANGpl{
\Question{
Poza rozwiązaniem równania  \( ax+b=cx+d \) algebraicznie, możesz przedstawić jego rozwiązanie za pomocą wykresu.  Kiedy linie  \( y=ax+b \) i \( y=cx+d \) są przedstawione graficznie, musisz znaleźć ich punkt przecięcia. Na rysunkach poniżej przedstawiono linie \( y=ax+b \) i \( y=cx+d\).  Wybierz rysunek, na którym równanie \( ax+b=cx+d \) ma tylko jedno nieujemne rozwiązanie.}
{}{}{}{}{}{}}
\LANGcs{
\Question{
Při grafickém řešení rovnice \( ax+b=cx+d \) postupujeme tak, že si nakreslíme grafy přímek \( y=ax+b \) a \( y=cx+d \) a hledáme jejich průsečík. Na kterém z obrázků jsou zakresleny přímky takové, že rovnice \( ax+b=cx+d \) má jediné nezáporné řešení?}
{}{}{}{}{}{}}
\LANGsk{
\Question{
Pri grafickom riešení rovnice \( ax+b=cx+d \) postupujeme tak,že si nakreslíme grafy priamok \( y=ax+b \) a \( y=cx+d \) a hľadáme ich priesečník. Na ktorom z obrázkov sú zakreslené priamky tak, že rovnica  \( ax+b=cx+d \) má jediné nezáporné riešenie?.}
{}{}{}{}{}{}}
\LANGes{
\Question{
NA}
{}{}{}{}{}{}}

\IMGanswer{1103049404a.pdf}
\IMGanswer{1103049404b.pdf}
\IMGanswer{1103049404c.pdf}
\IMGanswer{1103049404d.pdf}\End

\Data\ID{9000024101}
\Updated{2020-07-15 03:07:34}
\Level{A}
\SubArea{Linear equations and inequalities}
\LANGen{
\Question{
Identify the optimal first step to solve the following equation. The operation is
intended to be used on both sides of the equation. 
\[ 
 3x + 2 = -5x + 1
\]}
{add \((5x - 2)\)}{multiply by \(\frac{1}
{3}\)}{multiply by \(-\frac{1} 
{5}\)}{add \((-3x + 2)\)}{add \((5x + 1)\)}{add \((3x - 1)\)}}
\LANGpl{
\Question{
Które działanie należy wykonać jako pierwsze, aby rozwiązać podane równanie? Działanie to należy wykonać po obydwu stronach równania.  
\[ 
 3x + 2 = -5x + 1
\]}
{dodać \((5x - 2)\)}{pomnożyć przez \(\frac{1}
{3}\)}{pomnożyć przez \(-\frac{1} 
{5}\)}{dodać \((-3x + 2)\)}{dodać \((5x + 1)\)}{dodać \((3x - 1)\)}}
\LANGcs{
\Question{
Z nabídnutých možností vyberte nejvhodnější ekvivalentní úpravu,
pomocí které začneme řešit danou rovnici. Operace je zamýšlena pro aplikaci na obě strany rovnice.
\[ 
 3x + 2 = -5x + 1
\]}
{přičtení\( (5x-2) \)}{vynásobení číslem \(\frac{1}{3}\)}{vynásobení číslem \(-\frac{1} {5}\)}{přičtení \( (-3x+2) \)}{přičtení \( (5x + 1) \)}{přičtení \( (3x-1) \)}}
\LANGsk{
\Question{
Z ponúkaných možností vyberte najvhodnejšiu ekvivalentnú úpravu, pomocou ktorej začnete riešiť danú rovnicu. Táto úprava bude použitá na obidve strany rovnice. 
\[ 
 3x + 2 = -5x + 1
\]}
{pripočítanie \((5x - 2)\)}{vynásobenie číslom \(\frac{1}
{3}\)}{vynásobenie číslom \(-\frac{1} 
{5}\)}{pripočítanie \((-3x + 2)\)}{pripočítanie \((5x + 1)\)}{pripočítanie \((3x - 1)\)}}
\LANGes{
\Question{
Identifica el primer paso óptimo para resolver la siguiente ecuación. La operación será utilizada en ambos lados de la ecuación. 
\[ 
 3x + 2 = -5x + 1
\]}
{sumar  \((5x - 2)\)}{multiplicar por \(\frac{1}
{3}\)}{multiplicar por \(-\frac{1} 
{5}\)}{sumar  \((-3x + 2)\)}{sumar  \((5x + 1)\)}{sumar  \((3x - 1)\)}}
\End

\Data\ID{9000024102}
\Updated{2020-07-15 03:07:34}
\Level{A}
\SubArea{Linear equations and inequalities}
\LANGen{
\Question{
Identify the optimal first step to solve the following equation. The operation is
intended to be used on both sides of the equation. 
\[ 
 x + \frac{x} 
{6} = \frac{x} 
{15} + 1
\]}
{multiply by \(30\)}{multiply by \(6\)}{multiply by \(15\)}{subtract \((1 + x)\)}{subtract \(\left (\frac{x}
{6} + \frac{x} 
{15}\right )\)}{subtract \(\left (\frac{x}
{6} + 1\right )\)}}
\LANGpl{
\Question{
Które działanie należy wykonać jako pierwsze, aby rozwiązać podane równanie? Działanie to należy wykonać po obydwu stronach równania.  
\[ 
 x + \frac{x} 
{6} = \frac{x} 
{15} + 1
\]}
{pomnożyć przez \(30\)}{pomnożyć przez \(6\)}{pomnożyć przez  \(15\)}{odjąć  \((1 + x)\)}{odjąć  \(\left (\frac{x}
{6} + \frac{x} 
{15}\right )\)}{odjąć \(\left (\frac{x}
{6} + 1\right )\)}}
\LANGcs{
\Question{
Z nabídnutých možností vyberte nejvhodnější ekvivalentní úpravu,
pomocí které začneme řešit danou rovnici. Operace je zamýšlena pro aplikaci na obě strany rovnice.
\[ 
 x + \frac{x} 
{6} = \frac{x} 
{15} + 1
\]}
{vynásobení číslem \(30\)}{vynásobení číslem \(6\)}{vynásobení číslem \(15\)}{odečtení \( (1 + x) \)}{odečtení \( \left( \frac{x} {6} + \frac{x} {15}\right) \)}{odečtení \( \left( \frac{x} {6} + 1\right) \)}}
\LANGsk{
\Question{
Z ponúkaných možností vyberte najvhodnejšiu ekvivalentnú úpravu, pomocou ktorej začnete riešiť danú rovnicu. Táto úprava bude použitá na obidve strany rovnice. 
\[ 
 x + \frac{x} 
{6} = \frac{x} 
{15} + 1
\]}
{vynásobenie číslom \(30\)}{vynásobenie číslom \(6\)}{vynásobenie číslom \(15\)}{odpočítanie \((1 + x)\)}{odpočítanie \(\left (\frac{x}
{6} + \frac{x} 
{15}\right )\)}{odpočítanie \(\left (\frac{x}
{6} + 1\right )\)}}
\LANGes{
\Question{
Identifica el primer paso óptimo para resolver la siguiente ecuación. La operación será utilizada en ambos lados de la ecuación.
\[ 
 x + \frac{x} 
{6} = \frac{x} 
{15} + 1
\]}
{multiplicar por \(30\)}{multiplicar por \(6\)}{multiplicar por \(15\)}{restar  \((1 + x)\)}{restar  \(\left (\frac{x}
{6} + \frac{x} 
{15}\right )\)}{restar  \(\left (\frac{x}
{6} + 1\right )\)}}
\End

\Data\ID{9000024103}
\Updated{2020-07-15 03:07:34}
\Level{A}
\SubArea{Linear equations and inequalities}
\LANGen{
\Question{
Identify the optimal first step to solve the following equation. The operation is
intended to be used on both sides of the equation. 
\[ 
 \frac{x + 5} 
 {9} -\frac{x} 
{6} = \frac{x - 2} 
 {9} + \frac{x - 3} 
 {9}
\]}
{multiply by \(18\)}{multiply by \(6\)}{multiply by \(9\)}{multiply by \(54\)}{multiply by \(\frac{1}
{9}\)}{multiply by \(\frac{1}
{18}\)}}
\LANGpl{
\Question{
Które działanie należy wykonać jako pierwsze, aby rozwiązać podane równanie? Działanie to należy wykonać po obydwu stronach równania.  
\[ 
 \frac{x + 5} 
 {9} -\frac{x} 
{6} = \frac{x - 2} 
 {9} + \frac{x - 3} 
 {9}
\]}
{pomnożyć przez \(18\)}{pomnożyć przez \(6\)}{pomnożyć przez  \(9\)}{pomnożyć przez \(54\)}{pomnożyć przez \(\frac{1}
{9}\)}{pomnożyć przez  \(\frac{1}
{18}\)}}
\LANGcs{
\Question{
Z nabídnutých možností vyberte nejvhodnější ekvivalentní úpravu,
pomocí které začneme řešit danou rovnici. Operace je zamýšlena pro aplikaci na obě strany rovnice.
\[ 
 \frac{x + 5} 
 {9} -\frac{x} 
{6} = \frac{x - 2} 
 {9} + \frac{x - 3} 
 {9}
\]}
{vynásobení číslem \(18\)}{vynásobení číslem \(6\)}{vynásobení číslem \(9\)}{vynásobení číslem \(54\)}{vynásobení číslem \(\frac{1} {9}\)}{vynásobení číslem \(\frac{1} {18}\)}}
\LANGsk{
\Question{
Z ponúkaných možností vyberte najvhodnejšiu ekvivalentnú úpravu, pomocou ktorej začnete riešiť danú rovnicu. Úprava sa aplikuje na obidve strany rovnice. 
\[ 
 \frac{x + 5} 
 {9} -\frac{x} 
{6} = \frac{x - 2} 
 {9} + \frac{x - 3} 
 {9}
\]}
{vynásobenie číslom \(18\)}{vynásobenie číslom \(6\)}{vynásobenie číslom \(9\)}{vynásobenie číslom \(54\)}{vynásobenie číslom \(\frac{1}
{9}\)}{vynásobenie číslom \(\frac{1}
{18}\)}}
\LANGes{
\Question{
Identifica el primer paso óptimo para resolver la siguiente ecuación. La operación será utilizada en ambos lados de la ecuación.
\[ 
 \frac{x + 5} 
 {9} -\frac{x} 
{6} = \frac{x - 2} 
 {9} + \frac{x - 3} 
 {9}
\]}
{multiplicar por \(18\)}{multiplicar por \(6\)}{multiplicar por \(9\)}{multiplicar por \(54\)}{multiplicar por \(\frac{1}
{9}\)}{multiplicar por \(\frac{1}
{18}\)}}
\End

\Data\ID{9000024104}
\Updated{2020-07-15 03:07:34}
\Level{A}
\SubArea{Linear equations and inequalities}
\LANGen{
\Question{
Identify the optimal first step to solve the following equation. The operation is
intended to be used on both sides of the equation. 
\[ 
 5x = \frac{2 + x} 
 {5}
\]}
{multiply by \(5\)}{multiply by \(\frac{1}
{5}\)}{multiply by \(\frac{1}
{2}\)}{multiply by \(2\)}{multiply by \(\frac{1}
{x}\),
assuming \(x\neq 0\)}{multiply by \(x\),
assuming \(x\neq 0\)}}
\LANGpl{
\Question{
Które działanie należy wykonać jako pierwsze, aby rozwiązać podane równanie? Działanie to należy wykonać po obydwu stronach równania.  
\[ 
 5x = \frac{2 + x} 
 {5}
\]}
{pomnożyć przez  \(5\)}{pomnożyć przez \(\frac{1}
{5}\)}{pomnożyć przez \(\frac{1}
{2}\)}{pomnożyć przez  \(2\)}{pomnożyć przez \(\frac{1}
{x}\),
zakładając, że \(x\neq 0\)}{pomnożyć przez \(x\),
zakładając, że \(x\neq 0\)}}
\LANGcs{
\Question{
Z nabídnutých možností vyberte nejvhodnější ekvivalentní úpravu,
pomocí které začneme řešit danou rovnici. Operace je zamýšlena pro aplikaci na obě strany rovnice.
\[ 
 5x = \frac{2 + x} 
 {5}
\]}
{vynásobení číslem \(5\)}{vynásobení číslem \(\frac{1}{5}\)}{vynásobení číslem \(\frac{1}{2}\)}{vynásobení číslem \(2\)}{vynásobení výrazem \(\frac{1} {x}\) za předpokladu \(x\neq 0\)}{vynásobení výrazem \(x\) za předpokladu \(x\neq 0\)}}
\LANGsk{
\Question{
Z ponúkaných možností vyberte najvhodnejšiu ekvivalentnú úpravu, pomocou ktorej začnete riešiť danú rovnicu. Táto úprava bude použitá na obidve strany rovnice. 
\[ 
 5x = \frac{2 + x} 
 {5}
\]}
{vynásobenie číslom \(5\)}{vynásobenie číslom \(\frac{1}
{5}\)}{vynásobenie číslom \(\frac{1}
{2}\)}{vynásobenie číslom \(2\)}{vynásobenie výrazom \(\frac{1}
{x}\)
za predpokladu \(x\neq 0\)}{vynásobenie \(x\)
za predpokladu \(x\neq 0\)}}
\LANGes{
\Question{
Identifica el primer paso óptimo para resolver la siguiente ecuación. La operación será utilizada en ambos lados de la ecuación.
\[ 
 5x = \frac{2 + x} 
 {5}
\]}
{multiplicar por \(5\)}{multiplicar por \(\frac{1}
{5}\)}{multiplicar por \(\frac{1}
{2}\)}{multiplicar por \(2\)}{multiplicar por \(\frac{1}
{x}\),
suponiendo que \(x\neq 0\)}{multiplicar por \(x\),
suponiendo que \(x\neq 0\)}}
\End

\Data\ID{9000024107}
\Updated{2020-07-15 03:07:34}
\Level{A}
\SubArea{Linear equations and inequalities}
\LANGen{
\Question{
Identify the optimal first step to solve the following equation. The operation is
intended to be used on both sides of the equation. 
\[ 
 8x = \frac{x + 1} 
 {4} + 1
\]}
{multiply by \(4\)}{multiply by \(\frac{1}
{8}\)}{multiply by \(\frac{1}
{4}\)}{multiply by \((x + 1)\),
assuming \(x\neq - 1\)}{subtract \((x + 1)\)}{subtract \(1\)}}
\LANGpl{
\Question{
Które działanie należy wykonać jako pierwsze, aby rozwiązać podane równanie? Działanie to należy wykonać po obydwu stronach równania.  
\[ 
 8x = \frac{x + 1} 
 {4} + 1
\]}
{pomnożyć przez  \(4\)}{pomnożyć przez \(\frac{1}
{8}\)}{pomnożyć przez \(\frac{1}
{4}\)}{pomnożyć przez \((x + 1)\),
zakładając, że \(x\neq - 1\)}{odjąć \((x + 1)\)}{odjąć \(1\)}}
\LANGcs{
\Question{
Z nabídnutých možností vyberte nejvhodnější ekvivalentní úpravu,
pomocí které začneme řešit danou rovnici. Operace je zamýšlena pro aplikaci na obě strany rovnice.
\[ 
 8x = \frac{x + 1}{4} + 1
\]}
{vynásobení číslem \(4\)}{vynásobení číslem \(\frac{1} {8}\)}{vynásobení číslem \(\frac{1} {4}\)}{vynásobení výrazem \(x + 1\) za předpokladu \(x\neq -1\)}{odečtení \( (x + 1) \)}{odečtení \(1\)}}
\LANGsk{
\Question{
Z ponúkaných možností vyberte najvhodnejšiu ekvivalentnú úpravu, pomocou ktorej začnete riešiť danú rovnicu. Táto úprava bude použitá na obidve strany rovnice. 
\[ 
 8x = \frac{x + 1} 
 {4} + 1
\]}
{vynásobenie číslom \(4\)}{vynásobenie číslom \(\frac{1}
{8}\)}{vynásobenie číslom \(\frac{1}
{4}\)}{vynásobenie výrazom \((x + 1)\)
za predpokladu \(x\neq - 1\)}{odpočítanie \((x + 1)\)}{odpočítanie čísla \(1\)}}
\LANGes{
\Question{
Identifica el primer paso óptimo para resolver la siguiente ecuación. La operación será utilizada en ambos lados de la ecuación.
\[ 
 8x = \frac{x + 1} 
 {4} + 1
\]}
{multiplicar por \(4\)}{multiplicar por \(\frac{1}
{8}\)}{multiplicar por \(\frac{1}
{4}\)}{multiplicar por \((x + 1)\),
suponiendo que \(x\neq - 1\)}{restar  \((x + 1)\)}{restar  \(1\)}}
\End

\Data\ID{9000024108}
\Updated{2020-07-15 03:07:34}
\Level{A}
\SubArea{Linear equations and inequalities}
\LANGen{
\Question{
Identify the optimal first step to solve the following equation. The operation is
intended to be used on both sides of the equation. 
\[ 
 \frac{x + 1} 
 {2} -\frac{x - 2} 
 {3} = \frac{x} 
{4}
\]}
{multiply by \(12\)}{multiply by \(2\)}{multiply by \(3\)}{multiply by \(4\)}{multiply by \(24\)}{multiply by \((2x + 1)(x - 2)x\),
assuming \(x\not \in \left \{-\frac{1}
{2};2;0\right \}\)}}
\LANGpl{
\Question{
Które działanie należy wykonać jako pierwsze, aby rozwiązać podane równanie? Działanie to należy wykonać po obydwu stronach równania.  
\[ 
 \frac{x + 1} 
 {2} -\frac{x - 2} 
 {3} = \frac{x} 
{4}
\]}
{pomnożyć przez \(12\)}{pomnożyć przez \(2\)}{pomnożyć przez  \(3\)}{pomnożyć przez  \(4\)}{pomnożyć przez \(24\)}{pomnożyć przez \((2x + 1)(x - 2)x\),
zakładając, że \(x\not \in \left \{-\frac{1}
{2};2;0\right \}\)}}
\LANGcs{
\Question{
Z nabídnutých možností vyberte nejvhodnější ekvivalentní úpravu,
pomocí které začneme řešit danou rovnici. Operace je zamýšlena pro aplikaci na obě strany rovnice.
\[ 
 \frac{x + 1} 
 {2} -\frac{x - 2} 
 {3} = \frac{x} 
{4}
\]}
{vynásobení číslem \(12\)}{vynásobení číslem \(2\)}{vynásobení číslem \(3\)}{vynásobení číslem \(4\)}{vynásobení číslem \(24\)}{vynásobení výrazem \((2x + 1)(x - 2)x\) za předpokladu \(x\not \in \left \{-\frac{1} {2};2;0\right \}\)}}
\LANGsk{
\Question{
Z ponúkaných možností vyberte najvhodnejšiu ekvivalentnú úpravu, pomocou ktorej začnete riešiť danú rovnicu. Táto úprava bude použitá na obidve strany rovnice. 
\[ 
 \frac{x + 1} 
 {2} -\frac{x - 2} 
 {3} = \frac{x} 
{4}
\]}
{vynásobenie číslom \(12\)}{vynásobenie číslom \(2\)}{vynásobenie číslom \(3\)}{vynásobenie číslom \(4\)}{vynásobenie číslom \(24\)}{vynásobenie výrazom \((2x + 1)(x - 2)x\),
za predpokladu \(x\not \in \left \{-\frac{1}
{2};2;0\right \}\)}}
\LANGes{
\Question{
Identifica el primer paso óptimo para resolver la siguiente ecuación. La operación será utilizada en ambos lados de la ecuación.
\[ 
 \frac{x + 1} 
 {2} -\frac{x - 2} 
 {3} = \frac{x} 
{4}
\]}
{multiplicar por \(12\)}{multiplicar por \(2\)}{multiplicar por \(3\)}{multiplicar por \(4\)}{multiplicar por \(24\)}{multiplicar por \((2x + 1)(x - 2)x\),
suponiendo que \(x\not \in \left \{-\frac{1}
{2};2;0\right \}\)}}
\End

\Data\ID{9000024110}
\Updated{2020-07-15 03:07:34}
\Level{A}
\SubArea{Linear equations and inequalities}
\LANGen{
\Question{
Identify the optimal first step to solve the following equation. The operation is
intended to be used on both sides of the equation. 
\[ 
 11x - 2 = 2 - 4x
\]}
{add \((4x + 2)\)}{multiply by \(\frac{1}
{11}\)}{multiply by \(\left (-\frac{1}
{4}\right )\)}{add \((-11x + 4x)\)}{subtract \((4x + 2)\)}{add \((4x + 2)\)}}
\LANGpl{
\Question{
Które działanie należy wykonać jako pierwsze, aby rozwiązać podane równanie? Działanie to należy wykonać po obydwu stronach równania.  
\[ 
 11x - 2 = 2 - 4x
\]}
{dodać \((4x + 2)\)}{pomnożyć przez  \(\frac{1}
{11}\)}{pomnożyć przez  \(\left (-\frac{1}
{4}\right )\)}{dodać \((-11x + 4x)\)}{odjąć \((4x + 2)\)}{dodać \((4x + 2)\)}}
\LANGcs{
\Question{
Z nabídnutých možností vyberte nejvhodnější ekvivalentní úpravu,
pomocí které začneme řešit danou rovnici. Operace je zamýšlena pro aplikaci na obě strany rovnice.
\[ 
 11x - 2 = 2 - 4x
\]}
{přičtení \( (4x+2) \)}{vynásobení číslem \(\frac{1} {11}\)}{vynásobení číslem \(\left (-\frac{1} {4}\right )\)}{přičtení \( (- 11x+ 4x) \)}{odečtení \( (4x+ 2) \)}{přičtení \( (4x-2) \)}}
\LANGsk{
\Question{
Z ponúkaných možností vyberte najvhodnejšiu ekvivalentnú úpravu, pomocou ktorej začnete riešiť danú rovnicu. Táto úprava bude použitá na obidve strany rovnice. 
\[ 
 11x - 2 = 2 - 4x
\]}
{pripočítanie výrazu \((4x + 2)\)}{vynásobenie číslom \(\frac{1}
{11}\)}{vynásobenie číslom \(\left (-\frac{1}
{4}\right )\)}{pripočítanie \((-11x + 4x)\)}{odčítanie \((4x + 2)\)}{pripočítanie \((4x + 2)\)}}
\LANGes{
\Question{
Identifica el primer paso óptimo para resolver la siguiente ecuación. La operación será utilizada en ambos lados de la ecuación.
\[ 
 11x - 2 = 2 - 4x
\]}
{sumar  \((4x + 2)\)}{multiplicar por \(\frac{1}
{11}\)}{multiplicar por \(\left (-\frac{1}
{4}\right )\)}{sumar  \((-11x + 4x)\)}{restar  \((4x + 2)\)}{sumar  \((4x + 2)\)}}
\End

\Data\ID{1003031102}
\Updated{2020-07-15 03:07:12}
\Level{B}
\SubArea{Linear equations and inequalities}
\LANGen{
\Question{
The intended capacity of an IT course was filled by \( 20 \) people  interested to attend.
How many participants can check out without decreasing the course occupancy under \( 72\% \)?}
{at most \( 5 \)}{at most \( 6 \)}{at most \( 15 \)}{at most \( 14 \)}{}{}}
\LANGpl{
\Question{
Zamierzona ilość uczestników kursu informatycznego wyniosła
 \( 20 \) osób. Ilu uczestników może wypisać się bez zmniejszenia obłożenia kursu poniżej  \( 72\% \)?}
{najwyżej \( 5 \)}{najwyżej  \( 6 \)}{najwyżej \( 15 \)}{najwyżej  \( 14 \)}{}{}}
\LANGcs{
\Question{
\( 20 \) zájemců naplnilo plánovanou kapacitu počítačového kurzu. Kolik účastníků se může odhlásit, aniž by klesla obsazenost kurzu pod \( 72\,\% \)?}
{nejvýše \( 5 \)}{nejvýše \( 6 \)}{nejvýše \( 15 \)}{nejvýše \( 14 \)}{}{}}
\LANGsk{
\Question{
Plánovanú kapacitu počítačového kurzu naplnilo \( 20 \) záujemcov.
Koľko účastníkov sa môže odhlásiť, aby obsadenosť kurzu neklesla pod \( 72\% \)?}
{najviac \( 5 \)}{najviac \( 6 \)}{najviac \( 15 \)}{najviac \( 14 \)}{}{}}
\LANGes{
\Question{
La capacidad prevista de un curso de Informática fue ocupada por \( 20 \) personas interesadas en asistir. ¿Cuántos participantes pueden retirarse sin disminuir la ocupación del curso por debajo del \( 72\% \)?}
{a lo sumo \( 5 \)}{a lo sumo \( 6 \)}{a lo sumo \( 15 \)}{a lo sumo \( 14 \)}{}{}}
\End

\Data\ID{1003037501}
\Updated{2020-07-15 03:07:12}
\Level{B}
\SubArea{Linear equations and inequalities}
\LANGen{
\Question{
The difference of one-fifth of the unknown number and one-half of the number is greater than the difference of the number and thirteen. How many natural numbers do satisfy the given condition? (Note: Natural numbers are corresponding to positive integers.)}
{\( 9 \)}{\( 10 \)}{infinity}{\( 18 \)}{}{}}
\LANGpl{
\Question{
Różnica jednej piątej niewiadomej liczby i jednej drugiej tej liczby jest wyższa niż różnica tej liczby i trzynastu. Ile liczb spełnia dany warunek?}
{\( 9 \)}{\( 10 \)}{nieskończenie wiele}{\( 18 \)}{}{}}
\LANGcs{
\Question{
Rozdíl pětiny neznámého čísla a poloviny neznámého čísla je větší než rozdíl neznámého čísla a třinácti. Kolik přirozených čísel vyhovuje dané podmínce?}
{\( 9 \)}{\( 10 \)}{nekonečně mnoho}{\( 18 \)}{}{}}
\LANGsk{
\Question{
Rozdiel pätiny neznámeho čísla a polovice neznámeho čísla je väčší než rozdiel neznámeho čísla a trinásť. Koľko prirodzených čísel vyhovuje danej podmienke?}
{\( 9 \)}{\( 10 \)}{nekonečne veľa}{\( 18 \)}{}{}}
\LANGes{
\Question{
La diferencia entre la quinta parte de un número desconocido y la mitad del número es mayor que la diferencia entre el número y trece. ¿Cuántos números naturales cumplen la condición dada?}
{\( 9 \)}{\( 10 \)}{infinitos}{\( 18 \)}{}{}}
\End

\Data\ID{1003037502}
\Updated{2020-07-15 03:07:12}
\Level{B}
\SubArea{Linear equations and inequalities}
\LANGen{
\Question{
Five more than one-third of the unknown number is smaller than three times the difference of this number and five. Which numbers do satisfy the given condition?}
{greater than \( \frac{15}2 \)}{greater than \( \frac{15}4 \)}{greater than \( \frac{25}4 \)}{greater than \( \frac52 \)}{}{}}
\LANGpl{
\Question{
Pięć więcej niż jedna trzecia niewiadomej liczby jest mniejsze niż trzykrotność różnicy tej liczby i pięciu. Które liczby spełniają dany warunek?}
{więcej niż \( \frac{15}2 \)}{więcej niż \( \frac{15}4 \)}{więcej niż \( \frac{25}4 \)}{więcej niż \( \frac52 \)}{}{}}
\LANGcs{
\Question{
Třetina neznámého čísla zvětšená o pět je menší než trojnásobek rozdílu neznámého čísla a pěti. Jaká čísla splňují uvedenou podmínku?}
{větší než \( \frac{15}2 \)}{větší než \( \frac{15}4 \)}{větší než \( \frac{25}4 \)}{větší než \( \frac52 \)}{}{}}
\LANGsk{
\Question{
Tretina neznámeho čísla zväčšená o päť  je menšia než trojnásobok rozdielu neznámeho čísla a päť . Aké čísla vyhovujú danej podmienke?}
{väčšie než \( \frac{15}2 \)}{väčšie než \( \frac{15}4 \)}{väčšie než \( \frac{25}4 \)}{väčšie než \( \frac52 \)}{}{}}
\LANGes{
\Question{
La tercera parte de un número desconocido aumentado en cinco es menor que tres veces la diferencia entre ese número y cinco. ¿Qué números cumplen la condición?}
{mayores que \( \frac{15}2 \)}{mayores que \( \frac{15}4 \)}{mayores que \( \frac{25}4 \)}{mayores que \( \frac52 \)}{}{}}
\End

\Data\ID{1003046901}
\Updated{2020-07-15 03:07:12}
\Level{B}
\SubArea{Linear equations and inequalities}
\LANGen{
\Question{
We are given the inequality  \( -2x-\frac52 > 5-\frac x3 \).  Decide which of the following inequalities is equivalent to the given inequality, i.e. which of the following inequalities was obtained from the given inequality by equivalent transformations.}
{\( 12x+15 < 2x-30 \)}{\( 12x+15 > 2x-30 \)}{\( 12x-15 > 2x-30 \)}{\( 12x-15 < 2x+30 \)}{}{}}
\LANGpl{
\Question{
Dana jest nierówność  \( -2x-\frac52 > 5-\frac x3 \).  Zdecyduj, która z poniższych nierówności jest równoważna danej nierówności, tzn. która z poniższych nierówności powstała z danej nierówności w wyniki takich samych przekształceń.}
{\( 12x+15 < 2x-30 \)}{\( 12x+15 > 2x-30 \)}{\( 12x-15 > 2x-30 \)}{\( 12x-15 < 2x+30 \)}{}{}}
\LANGcs{
\Question{
Je dána nerovnice   \( -2x-\frac52 > 5-\frac x3 \).  Která z následujících nerovnic je s ní ekvivalentní, tj. vznikla ekvivalentní úpravou zadané nerovnice?}
{\( 12x+15 < 2x-30 \)}{\( 12x+15 > 2x-30 \)}{\( 12x-15 > 2x-30 \)}{\( 12x-15 < 2x+30 \)}{}{}}
\LANGsk{
\Question{
Je daná nerovnica  \( -2x-\frac52 > 5-\frac x3 \).  Ktorá z nasledujúcich nerovníc je s ňou ekvivalentná, t.j. vznikla ekvivalentnými úpravami danej nerovnice?}
{\( 12x+15 < 2x-30 \)}{\( 12x+15 > 2x-30 \)}{\( 12x-15 > 2x-30 \)}{\( 12x-15 < 2x+30 \)}{}{}}
\LANGes{
\Question{
Dada la inecuación  \( -2x-\frac52 > 5-\frac x3 \).  Determina cuál de las siguientes inecuaciones es equivalente a la inecuación dada.}
{\( 12x+15 < 2x-30 \)}{\( 12x+15 > 2x-30 \)}{\( 12x-15 > 2x-30 \)}{\( 12x-15 < 2x+30 \)}{}{}}
\End

\Data\ID{1003046902}
\Updated{2020-07-15 03:07:12}
\Level{B}
\SubArea{Linear equations and inequalities}
\LANGen{
\Question{
We are given the inequality \( 5-\frac{x+2}3 \leq \frac{2-x}6 \).  Decide which of the following inequalities is equivalent to the given inequality, i.e. which of the following inequalities was obtained from the given inequality by equivalent transformations.}
{\( 26-2x \leq 2-x \)}{\( 34-2x \leq 2-x \)}{\( 26-2x \geq 2-x \)}{\( 28-2x \leq 2-x \)}{}{}}
\LANGpl{
\Question{
Dana jest nierówność  \( 5-\frac{x+2}3 \leq \frac{2-x}6 \).  Zdecyduj, która z poniższych nierówności jest równoważna danej nierówności, tzn. która z poniższych nierówności powstała z danej nierówności w wyniki takich samych przekształceń.}
{\( 26-2x \leq 2-x \)}{\( 34-2x \leq 2-x \)}{\( 26-2x \geq 2-x \)}{\( 28-2x \leq 2-x \)}{}{}}
\LANGcs{
\Question{
Je dána nerovnice  \( 5-\frac{x+2}3 \leq \frac{2-x}6 \). Která z následujících nerovnic je s ní ekvivalentní, tj. vznikla ekvivalentní úpravou zadané nerovnice?}
{\( 26-2x \leq 2-x \)}{\( 34-2x \leq 2-x \)}{\( 26-2x \geq 2-x \)}{\( 28-2x \leq 2-x \)}{}{}}
\LANGsk{
\Question{
Je daná nerovnica \( 5-\frac{x+2}3 \leq \frac{2-x}6 \).  Ktorá z nasledujúcich nerovníc je s ňou ekvivalentná, t.j. vznikla ekvivalentnými úpravami zadanej nerovnice?}
{\( 26-2x \leq 2-x \)}{\( 34-2x \leq 2-x \)}{\( 26-2x \geq 2-x \)}{\( 28-2x \leq 2-x \)}{}{}}
\LANGes{
\Question{
Dada la inecuación \( 5-\frac{x+2}3 \leq \frac{2-x}6 \).  Determina cuál de las siguientes inecuaciones es equivalente a la inecuación dada.}
{\( 26-2x \leq 2-x \)}{\( 34-2x \leq 2-x \)}{\( 26-2x \geq 2-x \)}{\( 28-2x \leq 2-x \)}{}{}}
\End

\Data\ID{1003046903}
\Updated{2020-07-15 03:07:12}
\Level{B}
\SubArea{Linear equations and inequalities}
\LANGen{
\Question{
We are given the inequality  \( 3x+\frac{5-6x}2 > -2 \).  Decide which of the following inequalities is equivalent to the given inequality, i.e. which of the following inequalities was obtained from the given inequality by equivalent transformations.}
{\( 0\cdot x > -9 \)}{\( 0\cdot x > 9 \)}{\( 0\cdot x < -9 \)}{\( 3\cdot x > -9 \)}{}{}}
\LANGpl{
\Question{
Dana jest nierówność \( 3x+\frac{5-6x}2 > -2 \). Zdecyduj, która z poniższych nierówności jest równoważna danej nierówności, tzn. która z poniższych nierówności powstała z danej nierówności  w wyniki takich samych przekształceń.}
{\( 0\cdot x > -9 \)}{\( 0\cdot x > 9 \)}{\( 0\cdot x < -9 \)}{\( 3\cdot x > -9 \)}{}{}}
\LANGcs{
\Question{
Je dána nerovnice   \( 3x+\frac{5-6x}2 > -2 \).  Která z následujících nerovnic je s ní ekvivalentní, tj. vznikla ekvivalentní úpravou zadané nerovnice?}
{\( 0\cdot x > -9 \)}{\( 0\cdot x > 9 \)}{\( 0\cdot x < -9 \)}{\( 3\cdot x > -9 \)}{}{}}
\LANGsk{
\Question{
Je daná nerovnica  \( 3x+\frac{5-6x}2 > -2 \).  Ktorá z nasledujúcich nerovníc je s ňou ekvivalentná, t.j. vznikla ekvivalentnými úpravami zadanej nerovnice?}
{\( 0\cdot x > -9 \)}{\( 0\cdot x > 9 \)}{\( 0\cdot x < -9 \)}{\( 3\cdot x > -9 \)}{}{}}
\LANGes{
\Question{
Dada la inecuación  \( 3x+\frac{5-6x}2 > -2 \).  Determina cuál de las siguientes inecuaciones es equivalente a la inecuación dada.}
{\( 0\cdot x > -9 \)}{\( 0\cdot x > 9 \)}{\( 0\cdot x < -9 \)}{\( 3\cdot x > -9 \)}{}{}}
\End

\Data\ID{1003046904}
\Updated{2020-07-15 03:07:12}
\Level{B}
\SubArea{Linear equations and inequalities}
\LANGen{
\Question{
We are given the inequality  \( \frac1x+1 > \frac3{2x} \).  Decide which of the following inequalities has the different solution set than the given inequality has, i.e. choose the inequality which is not equivalent to the given inequality.}
{\( 1+x > \frac32 \)}{\( \frac2x+2>\frac3x \)}{\( 1>\frac3{2x}-\frac1x \)}{\( \frac1x-\frac3{2x}>-1 \)}{}{}}
\LANGpl{
\Question{
Dana jest nierówność  \( \frac1x+1 > \frac3{2x} \).  Zdecyduj, która z poniższych nierówności posiada inny zbiór rozwiązań niż podana nierówność, tzn. wybierz tę nierówność, która nie jest równoważna podanej wcześniej nierówności.}
{\( 1+x > \frac32 \)}{\( \frac2x+2>\frac3x \)}{\( 1>\frac3{2x}-\frac1x \)}{\( \frac1x-\frac3{2x}>-1 \)}{}{}}
\LANGcs{
\Question{
Je dána nerovnice   \( \frac1x+1 > \frac3{2x} \).  Z následujících nerovnic vyberte tu, která má jinou množinu kořenů než zadaná nerovnice, tj. není s ní ekvivalentní.}
{\( 1+x > \frac32 \)}{\( \frac2x+2>\frac3x \)}{\( 1>\frac3{2x}-\frac1x \)}{\( \frac1x-\frac3{2x}>-1 \)}{}{}}
\LANGsk{
\Question{
Je daná nerovnica  \( \frac1x+1 > \frac3{2x} \).  Z nasledujúcich nerovníc vyberte tú, ktorá má inú množinu koreňov, než zadaná nerovnica,  t.j. nie je s ňou ekvivalentná.}
{\( 1+x > \frac32 \)}{\( \frac2x+2>\frac3x \)}{\( 1>\frac3{2x}-\frac1x \)}{\( \frac1x-\frac3{2x}>-1 \)}{}{}}
\LANGes{
\Question{
Dada la inecuación  \( \frac1x+1 > \frac3{2x} \).  Determina cuál de las siguientes inecuaciones tiene un conjunto de raíces (soluciones) diferentes que la inecuación dada, es decir, no es equivalente a ella.}
{\( 1+x > \frac32 \)}{\( \frac2x+2>\frac3x \)}{\( 1>\frac3{2x}-\frac1x \)}{\( \frac1x-\frac3{2x}>-1 \)}{}{}}
\End

\Data\ID{1103049501}
\Updated{2020-07-15 03:07:12}
\Level{B}
\SubArea{Linear equations and inequalities}
\IMG{1103049501.pdf}
\LANGen{
\Question{
Find the linear inequality of which the solution set is shown in red in the graph below.}
{\( x+2\geq\frac{2-x}2 \)}{\( x+2\leq-\frac x2+1 \)}{\( x+2>\frac{2-x}2 \)}{\( x+2 < -\frac x2+1 \)}{}{}}
\LANGpl{
\Question{
Wybierz nierówność liniową, której zbiór rozwiązań oznaczono na poniższym wykresie kolorem czerwonym?}
{\( x+2\geq\frac{2-x}2 \)}{\( x+2\leq-\frac x2+1 \)}{\( x+2>\frac{2-x}2 \)}{\( x+2 < -\frac x2+1 \)}{}{}}
\LANGcs{
\Question{
Na obrázku jsou červeně vyznačena všechna řešení lineární nerovnice. Vyberte nerovnici, jejíž grafické řešení je na obrázku.}
{\( x+2\geq\frac{2-x}2 \)}{\( x+2\leq-\frac x2+1 \)}{\( x+2>\frac{2-x}2 \)}{\( x+2 < -\frac x2+1 \)}{}{}}
\LANGsk{
\Question{
Vyberte lineárnu nerovnicu, ktorej grafické riešenie je na obrázku vyznačené červenou farbou.}
{\( x+2\geq\frac{2-x}2 \)}{\( x+2\leq-\frac x2+1 \)}{\( x+2>\frac{2-x}2 \)}{\( x+2 < -\frac x2+1 \)}{}{}}
\LANGes{
\Question{
Determina la inecuación lineal cuyo conjunto de soluciones está representado por una semirecta roja en la imagen.}
{\( x+2\geq\frac{2-x}2 \)}{\( x+2\leq-\frac x2+1 \)}{\( x+2>\frac{2-x}2 \)}{\( x+2 < -\frac x2+1 \)}{}{}}
\End

\Data\ID{1103049502}
\Updated{2020-07-15 03:07:12}
\Level{B}
\SubArea{Linear equations and inequalities}
\IMG{1103049502.pdf}
\LANGen{
\Question{
Find the linear inequality of which the solution set is shown in red in the graph below.}
{\( 2x+2 < \frac x4-\frac 12 \)}{\( 2x+2 > \frac{x-2}4 \)}{\( 2x+2 < -\frac x2+2 \)}{\( \frac x2+2 < \frac x4-\frac12 \)}{}{}}
\LANGpl{
\Question{
Wybierz nierówność liniową, której zbiór rozwiązań oznaczono na poniższym wykresie kolorem czerwonym?}
{\( 2x+2 < \frac x4-\frac 12 \)}{\( 2x+2 > \frac{x-2}4 \)}{\( 2x+2 < -\frac x2+2 \)}{\( \frac x2+2 < \frac x4-\frac12 \)}{}{}}
\LANGcs{
\Question{
Na obrázku jsou červeně vyznačena všechna řešení lineární nerovnice. Vyberte nerovnici, jejíž grafické řešení je na obrázku.}
{\( 2x+2 < \frac x4-\frac 12 \)}{\( 2x+2 > \frac{x-2}4 \)}{\( 2x+2 < -\frac x2+2 \)}{\( \frac x2+2 < \frac x4-\frac12 \)}{}{}}
\LANGsk{
\Question{
Vyberte lineárnu nerovnicu, ktorej grafické riešenie je na obrázku vyznačené červenou farbou.}
{\( 2x+2 < \frac x4-\frac 12 \)}{\( 2x+2 > \frac{x-2}4 \)}{\( 2x+2 < -\frac x2+2 \)}{\( \frac x2+2 < \frac x4-\frac12 \)}{}{}}
\LANGes{
\Question{
Determina la inecuación lineal cuyo conjunto de soluciones está representado por una semirecta roja en la imagen.}
{\( 2x+2 < \frac x4-\frac 12 \)}{\( 2x+2 > \frac{x-2}4 \)}{\( 2x+2 < -\frac x2+2 \)}{\( \frac x2+2 < \frac x4-\frac12 \)}{}{}}
\End

\Data\ID{1103049503}
\Updated{2020-07-15 03:07:10}
\Level{B}
\SubArea{Linear equations and inequalities}
\LANGen{
\Question{
Find the picture where marked in red are all the solutions of the following inequality. 
\[ 2x-\frac{1+4x}4\geq-2 \]}
{}{}{}{}{}{}}
\LANGpl{
\Question{
Wybierz rysunek, w którym na czerwono przedstawiono wszystkie rozwiązania następującej nierówności. 
\[ 2x-\frac{1+4x}4\geq-2 \]}
{}{}{}{}{}{}}
\LANGcs{
\Question{
Vyberte obrázek, na kterém jsou červeně vyznačena všechna řešení následující nerovnice.
\[ 2x-\frac{1+4x}4\geq-2 \]}
{}{}{}{}{}{}}
\LANGsk{
\Question{
Vyberte obrázok, na ktorom sú červenou vyznačené všetky riešenia nasledujúcej nerovnice. 
\[ 2x-\frac{1+4x}4\geq-2 \]}
{}{}{}{}{}{}}
\LANGes{
\Question{
NA}
{}{}{}{}{}{}}

\IMGanswer{1103049503a.pdf}
\IMGanswer{1103049503b.pdf}
\IMGanswer{1103049503c.pdf}
\IMGanswer{1103049503d.pdf}\End

\Data\ID{1103049504}
\Updated{2020-07-15 03:07:12}
\Level{B}
\SubArea{Linear equations and inequalities}
\IMG{1103049504.pdf}
\LANGen{
\Question{
Which inequality has the solution set graphed in red in the picture?}
{\( \frac{5-2x}2 < \frac 72 \)}{\( \frac{5-2x}2 > \frac72 \)}{\( \frac{5-2x}2 > 4 \)}{\( \frac{5-2x}2 < 7 \)}{}{}}
\LANGpl{
\Question{
Wybierz nierówność, której zbiór rozwiązań oznaczono na poniższym wykresie kolorem czerwonym?}
{\( \frac{5-2x}2 < \frac 72 \)}{\( \frac{5-2x}2 > \frac72 \)}{\( \frac{5-2x}2 > 4 \)}{\( \frac{5-2x}2 < 7 \)}{}{}}
\LANGcs{
\Question{
Která nerovnice má množinu řešení, která je graficky znázorněna na obrázku?}
{\( \frac{5-2x}2 < \frac 72 \)}{\( \frac{5-2x}2 > \frac72 \)}{\( \frac{5-2x}2 > 4 \)}{\( \frac{5-2x}2 < 7 \)}{}{}}
\LANGsk{
\Question{
Ktorá nerovnica má množinu riešení, ktorá je graficky znázornená na obrázku?}
{\( \frac{5-2x}2 < \frac 72 \)}{\( \frac{5-2x}2 > \frac72 \)}{\( \frac{5-2x}2 > 4 \)}{\( \frac{5-2x}2 < 7 \)}{}{}}
\LANGes{
\Question{
¿Qué inecuación tiene el conjunto de soluciones representado gráficamente en la imagen?}
{\( \frac{5-2x}2 < \frac 72 \)}{\( \frac{5-2x}2 > \frac72 \)}{\( \frac{5-2x}2 > 4 \)}{\( \frac{5-2x}2 < 7 \)}{}{}}
\End

\Data\ID{1103049505}
\Updated{2020-07-15 03:07:10}
\Level{B}
\SubArea{Linear equations and inequalities}
\LANGen{
\Question{
Which graph shows in red the solution set of the following inequality?
\[ 2x \geq \frac{6-4x}2 \]}
{}{}{}{}{}{}}
\LANGpl{
\Question{
Który wykres przedstawia na czerwono  zbiór rozwiązań następującej nierówności? 
\[ 2x \geq \frac{6-4x}2 \]}
{}{}{}{}{}{}}
\LANGcs{
\Question{
Vyberte obrázek, na kterém jsou červeně vyznačena všechna řešení nerovnice.
\[ 2x\geq \frac{6-4x}2 \]}
{}{}{}{}{}{}}
\LANGsk{
\Question{
Vyberte obrázok, na ktorom sú červenou vyznačené všetky riešenia danej nerovnice.
\[ 2x \geq \frac{6-4x}2 \]}
{}{}{}{}{}{}}
\LANGes{
\Question{
NA}
{}{}{}{}{}{}}

\IMGanswer{1103049505a.pdf}
\IMGanswer{1103049505b.pdf}
\IMGanswer{1103049505c.pdf}
\IMGanswer{1103049505d.pdf}\End

\Data\ID{1103049506}
\Updated{2020-07-15 03:07:10}
\Level{B}
\SubArea{Linear equations and inequalities}
\LANGen{
\Question{
Which graph shows in red the solution set of the following inequality? (Note: If the graph does not show in red any interval, then it means that the inequality has no solution.)
\[ \frac x3-\frac{2x+1}6 < 4 \]}
{}{}{}{}{}{}}
\LANGpl{
\Question{
Który wykres przedstawia na czerwono zbiór rozwiązań następujących nierówności? (Uwaga: jeśli wykres nie pokazuje na czerwono żadnego przedziału, oznacza to, że nierówność  ta nie ma rozwiązania.)
\[ \frac x3-\frac{2x+1}6 < 4 \]}
{}{}{}{}{}{}}
\LANGcs{
\Question{
Vyberte obrázek, na němž je červeně vyznačena množina řešení následující nerovnice. (Poznámka: Pokud v obrázku není červeně vyznačena žádná množina, znamená to, že nerovnice nemá řešení.)
\[ \frac x3-\frac{2x+1}6 < 4 \]}
{}{}{}{}{}{}}
\LANGsk{
\Question{
Vyberte graf, na ktorom je červenou vyznačená množina riešení nasledujúcej nerovnice. (Poznámka: Ak na grafe nie je červenou vyznačená žiadna množina, znamená to, že nerovnica nemá riešenie.)
\[ \frac x3-\frac{2x+1}6 < 4 \]}
{}{}{}{}{}{}}
\LANGes{
\Question{
NA}
{}{}{}{}{}{}}

\IMGanswer{1103049506a.pdf}
\IMGanswer{1103049506b.pdf}
\IMGanswer{1103049506c.pdf}
\IMGanswer{1103049506d.pdf}\End

\Data\ID{9000018001}
\Updated{2020-07-15 03:07:34}
\Level{B}
\SubArea{Linear equations and inequalities}
\LANGen{
\Question{
Solve the following inequality. 
\[ 
 -3x > 6
\]}
{\(x\in \left (-\infty ;-2\right )\)}{\(x\in \left (-\infty ;-2\right ] \)}{\(x\in \left (-2;\infty \right )\)}{\(x\in \left [ -2;\infty \right )\)}{}{}}
\LANGpl{
\Question{
Rozwiąż następującą nierówność.
\[ 
 -3x > 6
\]}
{\(x\in \left (-\infty ;-2\right )\)}{\(x\in \left (-\infty ;-2\right ] \)}{\(x\in \left (-2;\infty \right )\)}{\(x\in \left [ -2;\infty \right )\)}{}{}}
\LANGcs{
\Question{
Vyřešte následující nerovnici.
\[- 3x > 6\]}
{\(\left (-\infty ;-2\right )\)}{\(\left (-\infty ;-2\right \rangle \)}{\(\left (-2;\infty \right )\)}{\(\left \langle -2;\infty \right )\)}{}{}}
\LANGsk{
\Question{
Vyriešte nasledujúcu nerovnicu.
\[ 
 -3x > 6
\]}
{\(x\in \left (-\infty ;-2\right )\)}{\(x\in \left (-\infty ;-2\right ] \)}{\(x\in \left (-2;\infty \right )\)}{\(x\in \left [ -2;\infty \right )\)}{}{}}
\LANGes{
\Question{
Resuelve la siguiente inecuación.
\[ 
 -3x > 6
\]}
{\(x\in \left (-\infty ;-2\right )\)}{\(x\in \left (-\infty ;-2\right ] \)}{\(x\in \left (-2;\infty \right )\)}{\(x\in \left [ -2;\infty \right )\)}{}{}}
\End

\Data\ID{9000018002}
\Updated{2020-07-15 03:07:34}
\Level{B}
\SubArea{Linear equations and inequalities}
\LANGen{
\Question{
Assuming positive integer \(x\),
solve the following inequality. 
\[ 
 -5x\geq - 1
\]}
{\(x\in \emptyset \)}{\(x\in \left \{0\right \}\)}{\(x\in \left (0; \frac{1} 
{5}\right ] \)}{\(x\in \left \{\frac{1}
{5}\right \}\)}{}{}}
\LANGpl{
\Question{
Zakładając, że \(x\) jest liczbą całkowitą dodatnią,
rozwiąż podaną nierówność.
\[ 
 -5x\geq - 1
\]}
{\(x\in \emptyset \)}{\(x\in \left \{0\right \}\)}{\(x\in \left (0; \frac{1} 
{5}\right ] \)}{\(x\in \left \{\frac{1}
{5}\right \}\)}{}{}}
\LANGcs{
\Question{
Určete množinu všech řešení dané nerovnice v oboru přirozených čísel.
\[- 5x\geq - 1\]}
{\(\emptyset \)}{\(\left \{0\right \}\)}{\(\left (0; \frac{1} 
{5}\right \rangle \)}{\(\left \{\frac{1}
{5}\right \}\)}{}{}}
\LANGsk{
\Question{
Vyriešte nerovnicu v obore prirodzených čísel.
\[ 
 -5x\geq - 1
\]}
{\(x\in \emptyset \)}{\(x\in \left \{0\right \}\)}{\(x\in \left (0; \frac{1} 
{5}\right ] \)}{\(x\in \left \{\frac{1}
{5}\right \}\)}{}{}}
\LANGes{
\Question{
Suponiendo que \(x\) es un número natural, resuelve la siguiente inecuación. 
\[ 
 -5x\geq - 1
\]}
{\(x\in \emptyset \)}{\(x\in \left \{0\right \}\)}{\(x\in \left (0; \frac{1} 
{5}\right ] \)}{\(x\in \left \{\frac{1}
{5}\right \}\)}{}{}}
\End

\Data\ID{9000018003}
\Updated{2020-07-15 03:07:34}
\Level{B}
\SubArea{Linear equations and inequalities}
\LANGen{
\Question{
Assuming \(x\in \left (0;3\right ] \),
solve the following inequality. 
\[ 
 6 - 2x\leq 3x - 4
\]}
{\(x\in \left [ 2;3\right ] \)}{\(x\in \left (0;3\right ] \)}{\(x\in \left (0;2\right ] \)}{\(x\in \left (0;\infty \right )\)}{}{}}
\LANGpl{
\Question{
Zakładając, że \(x\in \left (0;3\right ] \),
rozwiąż podaną nierówność.
\[ 
 6 - 2x\leq 3x - 4
\]}
{\(x\in \left [ 2;3\right ] \)}{\(x\in \left (0;3\right ] \)}{\(x\in \left (0;2\right ] \)}{\(x\in \left (0;\infty \right )\)}{}{}}
\LANGcs{
\Question{
Určete množinu všech řešení dané nerovnice v intervalu \(\left (0;3\right \rangle \). \[6 - 2x\leq 3x - 4\]}
{\(\left \langle 2;3\right \rangle \)}{\(\left (0;3\right \rangle \)}{\(\left (0;2\right \rangle \)}{\(\left (0;\infty \right )\)}{}{}}
\LANGsk{
\Question{
V intervale \(x\in \left (0;3\right ] \)
vyriešte nasledujúcu nerovnicu.
\[ 
 6 - 2x\leq 3x - 4
\]}
{\(x\in \left [ 2;3\right ] \)}{\(x\in \left (0;3\right ] \)}{\(x\in \left (0;2\right ] \)}{\(x\in \left (0;\infty \right )\)}{}{}}
\LANGes{
\Question{
Suponiendo que \(x\in \left (0;3\right ] \),
resuelve la siguiente inecuación. 
\[ 
 6 - 2x\leq 3x - 4
\]}
{\(x\in \left [ 2;3\right ] \)}{\(x\in \left (0;3\right ] \)}{\(x\in \left (0;2\right ] \)}{\(x\in \left (0;\infty \right )\)}{}{}}
\End

\Data\ID{9000018004}
\Updated{2020-07-15 03:07:34}
\Level{B}
\SubArea{Linear equations and inequalities}
\LANGen{
\Question{
Find the maximal integer which satisfies the following inequality:
\[ 
 2x - 5 < 4 - x
\]}
{\(2\)}{\(- 3\)}{\(- 2\)}{\(3\)}{}{}}
\LANGpl{
\Question{
Znajdź największą liczbę całkowitą, która spełnia podaną nierówność:
\[ 
 2x - 5 < 4 - x
\]}
{\(2\)}{\(- 3\)}{\(- 2\)}{\(3\)}{}{}}
\LANGcs{
\Question{
Najděte největší číslo \(x\in \mathbb{Z}\),
které je řešením nerovnice: \[2x - 5 < 4 - x\]}
{\(2\)}{\(- 3\)}{\(- 2\)}{\(3\)}{}{}}
\LANGsk{
\Question{
Nájdite najväčšie celé číslo, ktoré je riešením nerovnice:
\[ 
 2x - 5 < 4 - x
\]}
{\(2\)}{\(- 3\)}{\(- 2\)}{\(3\)}{}{}}
\LANGes{
\Question{
Encuentra el número mayor \(x\in \mathbb{Z}\),
que es solución de la siguiente inecuación: \[2x - 5 < 4 - x\]}
{\(2\)}{\(- 3\)}{\(- 2\)}{\(3\)}{}{}}
\End

\Data\ID{9000018006}
\Updated{2020-07-15 03:07:34}
\Level{B}
\SubArea{Linear equations and inequalities}
\LANGen{
\Question{
Assuming negative integer \(x\),
solve the following inequality:
\[ 
 x - 2 > 1 - x - 8
\]}
{\(x\in \left \{-2;-1\right \}\)}{\(x\in \left \{-3;-2;-1\right \}\)}{\(x\in \left \{-3;-2\right \}\)}{\(x\in \left \{-1\right \}\)}{}{}}
\LANGpl{
\Question{
Zakładając, że  \(x\) jest liczbą całkowitą ujemną, rozwiąż podaną nierówność:
\[ 
 x - 2 > 1 - x - 8
\]}
{\(x\in \left \{-2;-1\right \}\)}{\(x\in \left \{-3;-2;-1\right \}\)}{\(x\in \left \{-3;-2\right \}\)}{\(x\in \left \{-1\right \}\)}{}{}}
\LANGcs{
\Question{
Určete všechna záporná celá čísla, která jsou řešením nerovnice:
\[x - 2 > 1 - x - 8\]}
{\(\left \{-2;-1\right \}\)}{\(\left \{-3;-2;-1\right \}\)}{\(\left \{-3;-2\right \}\)}{\(\left \{-1\right \}\)}{}{}}
\LANGsk{
\Question{
Nájdite všetky záporné celé čísla, ktoré sú riešením nerovnice:
\[ 
 x - 2 > 1 - x - 8
\]}
{\(x\in \left \{-2;-1\right \}\)}{\(x\in \left \{-3;-2;-1\right \}\)}{\(x\in \left \{-3;-2\right \}\)}{\(x\in \left \{-1\right \}\)}{}{}}
\LANGes{
\Question{
Suponiendo que \(x\),
es un entero negativo, resuelve la siguiente inecuación:
\[ 
 x - 2 > 1 - x - 8
\]}
{\(x\in \left \{-2;-1\right \}\)}{\(x\in \left \{-3;-2;-1\right \}\)}{\(x\in \left \{-3;-2\right \}\)}{\(x\in \left \{-1\right \}\)}{}{}}
\End

\Data\ID{9000018101}
\Updated{2020-07-15 03:07:34}
\Level{B}
\SubArea{Linear equations and inequalities}
\LANGen{
\Question{
Solve the following inequality:
\[ 
 7 -\left (4x - 1\right ) < 3\left (x + 4\right )
\]}
{\(x\in \left (-\frac{4}
{7};\infty \right )\)}{\(x\in \left (-\infty ; \frac{4} 
{7}\right )\)}{\(x\in \left (\frac{4}
{7};\infty \right )\)}{\(x\in \left (-\infty ;-\frac{4}
{7}\right )\)}{}{}}
\LANGpl{
\Question{
Rozwiąż podaną nierówność: 
\[ 
 7 -\left (4x - 1\right ) < 3\left (x + 4\right )
\]}
{\(x\in \left (-\frac{4}
{7};\infty \right )\)}{\(x\in \left (-\infty ; \frac{4} 
{7}\right )\)}{\(x\in \left (\frac{4}
{7};\infty \right )\)}{\(x\in \left (-\infty ;-\frac{4}
{7}\right )\)}{}{}}
\LANGcs{
\Question{
Vyřešte danou nerovnici:
\[7 -\left (4x - 1\right ) < 3\left (x + 4\right )\]}
{\(x\in\left (-\frac{4}
{7};\infty \right )\)}{\(x\in\left (-\infty ; \frac{4} 
{7}\right )\)}{\(x\in\left (\frac{4}
{7};\infty \right )\)}{\(x\in\left (-\infty ;-\frac{4}
{7}\right )\)}{}{}}
\LANGsk{
\Question{
Vyriešte nasledujúcu nerovnicu:
\[ 
 7 -\left (4x - 1\right ) < 3\left (x + 4\right )
\]}
{\(x\in \left (-\frac{4}
{7};\infty \right )\)}{\(x\in \left (-\infty ; \frac{4} 
{7}\right )\)}{\(x\in \left (\frac{4}
{7};\infty \right )\)}{\(x\in \left (-\infty ;-\frac{4}
{7}\right )\)}{}{}}
\LANGes{
\Question{
Resuelve la siguiente inecuación:
\[ 
 7 -\left (4x - 1\right ) < 3\left (x + 4\right )
\]}
{\(x\in \left (-\frac{4}
{7};\infty \right )\)}{\(x\in \left (-\infty ; \frac{4} 
{7}\right )\)}{\(x\in \left (\frac{4}
{7};\infty \right )\)}{\(x\in \left (-\infty ;-\frac{4}
{7}\right )\)}{}{}}
\End

\Data\ID{9000018102}
\Updated{2020-07-15 03:07:34}
\Level{B}
\SubArea{Linear equations and inequalities}
\LANGen{
\Question{
Solve the following inequality:
\[ 
 \left (5 + 2x\right )\cdot \left (-3\right ) + 16 < 20 - 6x
\]}
{\(x\in \left (-\infty ;\infty \right )\)}{\(x\in \left (-\infty ;2\right )\)}{\(x\in \emptyset \)}{\(x\in \left (2;\infty \right )\)}{}{}}
\LANGpl{
\Question{
Rozwiąż podaną nierówność: 
\[ 
 \left (5 + 2x\right )\cdot \left (-3\right ) + 16 < 20 - 6x
\]}
{\(x\in \left (-\infty ;\infty \right )\)}{\(x\in \left (-\infty ;2\right )\)}{\(x\in \emptyset \)}{\(x\in \left (2;\infty \right )\)}{}{}}
\LANGcs{
\Question{
Vyřešte danou nerovnici:
\[\left (5 + 2x\right )\cdot \left (-3\right ) + 16 < 20 - 6x\]}
{\(x\in\left (-\infty ;\infty \right )\)}{\(x\in\left (-\infty ;2\right )\)}{\(x\in\emptyset \)}{\(x\in\left (2;\infty \right )\)}{}{}}
\LANGsk{
\Question{
Vyriešte nasledujúcu nerovnicu:
\[ 
 \left (5 + 2x\right )\cdot \left (-3\right ) + 16 < 20 - 6x
\]}
{\(x\in \left (-\infty ;\infty \right )\)}{\(x\in \left (-\infty ;2\right )\)}{\(x\in \emptyset \)}{\(x\in \left (2;\infty \right )\)}{}{}}
\LANGes{
\Question{
Resuelve la siguiente inecuación:
\[ 
 \left (5 + 2x\right )\cdot \left (-3\right ) + 16 < 20 - 6x
\]}
{\(x\in \left (-\infty ;\infty \right )\)}{\(x\in \left (-\infty ;2\right )\)}{\(x\in \emptyset \)}{\(x\in \left (2;\infty \right )\)}{}{}}
\End

\Data\ID{9000018103}
\Updated{2020-07-15 03:07:34}
\Level{B}
\SubArea{Linear equations and inequalities}
\LANGen{
\Question{
Solve the following inequality assuming positive integer
\(x\).
\[ 
 1\frac{1}
{3}\leq -\frac{x - 4} 
 {2}
\]}
{\(x\in \left \{1\right \}\)}{\(x\in \left \{0;1\right \}\)}{\(x\in \left (0; \frac{4} 
{3}\right ] \)}{\(x\in \emptyset \)}{}{}}
\LANGpl{
\Question{
Rozwiąż podaną nierówność, zakładając, że 
\(x\) jest dodatnią liczbą całkowitą.
\[ 
 1\frac{1}
{3}\leq -\frac{x - 4} 
 {2}
\]}
{\(x\in \left \{1\right \}\)}{\(x\in \left \{0;1\right \}\)}{\(x\in \left (0; \frac{4} 
{3}\right ] \)}{\(x\in \emptyset \)}{}{}}
\LANGcs{
\Question{
Vyřešte danou nerovnici v
oboru přirozených čísel.
\[1\frac{1}
{3}\leq -\frac{x-4} 
 {2} \]}
{\(x\in\left \{1\right \}\)}{\(x\in\left \{0;1\right \}\)}{\(x\in\left (0; \frac{4} 
{3}\right \rangle \)}{\(x\in\emptyset \)}{}{}}
\LANGsk{
\Question{
Vyriešte nasledujúcu nerovnicu v obore prirodzených čísel.
\[ 
 1\frac{1}
{3}\leq -\frac{x - 4} 
 {2}
\]}
{\(x\in \left \{1\right \}\)}{\(x\in \left \{0;1\right \}\)}{\(x\in \left (0; \frac{4} 
{3}\right ] \)}{\(x\in \emptyset \)}{}{}}
\LANGes{
\Question{
Suponiendo que \(x\) es un número natural, resuelve la siguiente inecuación
\[ 
 1\frac{1}
{3}\leq -\frac{x - 4} 
 {2}
\]}
{\(x\in \left \{1\right \}\)}{\(x\in \left \{0;1\right \}\)}{\(x\in \left (0; \frac{4} 
{3}\right ] \)}{\(x\in \emptyset \)}{}{}}
\End

\Data\ID{9000018104}
\Updated{2020-07-15 03:07:34}
\Level{B}
\SubArea{Linear equations and inequalities}
\LANGen{
\Question{
Find the maximal integer which solves the following inequality. 
\[ 
 1 - 3x > 3\left (4 - x\right ) + 2x
\]}
{\(- 6\)}{\(- 5\)}{\(- 3\)}{\(- 2\)}{}{}}
\LANGpl{
\Question{
Znajdź największą liczbę całkowitą, która jest rozwiązaniem następującej nierówności. 
\[ 
 1 - 3x > 3\left (4 - x\right ) + 2x
\]}
{\(- 6\)}{\(- 5\)}{\(- 3\)}{\(- 2\)}{}{}}
\LANGcs{
\Question{
Najděte největší číslo \(x\in \mathbb{Z}\),
které je řešením dané nerovnice. \[1 - 3x > 3\left (4 - x\right ) + 2x\]}
{\(- 6\)}{\(- 5\)}{\(- 3\)}{\(- 2\)}{}{}}
\LANGsk{
\Question{
Nájdite najväčšie celé číslo, ktoré je riešením danej nerovnice. 
\[ 
 1 - 3x > 3\left (4 - x\right ) + 2x
\]}
{\(- 6\)}{\(- 5\)}{\(- 3\)}{\(- 2\)}{}{}}
\LANGes{
\Question{
Encuentra el número mayor \(x\in \mathbb{Z}\),
que es la solución de la siguiente inecuación. \[1 - 3x > 3\left (4 - x\right ) + 2x\]}
{\(- 6\)}{\(- 5\)}{\(- 3\)}{\(- 2\)}{}{}}
\End

\Data\ID{9000018106}
\Updated{2020-07-15 03:07:34}
\Level{B}
\SubArea{Linear equations and inequalities}
\LANGen{
\Question{
Find the set of all the positive integers
\(x\) for which the
expression \(\frac{3x-7}
 {14} \) is
smaller than \(\frac{7-2x}
 {7} \).}
{\(\left \{1;2\right \}\)}{\(\left \{1;2;3;4\right \}\)}{\(\left \{1;2;3\right \}\)}{\(\left \{1\right \}\)}{}{}}
\LANGpl{
\Question{
Znajdź zbiór rozwiązań wszystkich liczb całkowitych dodatnich 
\(x\), dla których wyrażenie
 \(\frac{3x-7}
 {14} \) jest mniejsze niż \(\frac{7-2x}
 {7} \).}
{\(\left \{1;2\right \}\)}{\(\left \{1;2;3;4\right \}\)}{\(\left \{1;2;3\right \}\)}{\(\left \{1\right \}\)}{}{}}
\LANGcs{
\Question{
Pro která \(x\in \mathbb{N}\) je
zlomek \(\frac{3x-7}
 {14} \) menší
než zlomek \(\frac{7-2x}
 {7} \)?}
{\(\left \{1;2\right \}\)}{\(\left \{1;2;3;4\right \}\)}{\(\left \{1;2;3\right \}\)}{\(\left \{1\right \}\)}{}{}}
\LANGsk{
\Question{
Pre ktoré kladné celé čísla 
\(x\) je výraz
 \(\frac{3x-7}
 {14} \) 
menší ako výraz \(\frac{7-2x}
 {7} \)?}
{\(\left \{1;2\right \}\)}{\(\left \{1;2;3;4\right \}\)}{\(\left \{1;2;3\right \}\)}{\(\left \{1\right \}\)}{}{}}
\LANGes{
\Question{
¿Para qué \(x\in \mathbb{N}\)  la fracción\(\frac{3x-7}
 {14} \) es menor que la fracción \(\frac{7-2x}
 {7} \)?}
{\(\left \{1;2\right \}\)}{\(\left \{1;2;3;4\right \}\)}{\(\left \{1;2;3\right \}\)}{\(\left \{1\right \}\)}{}{}}
\End

\Data\ID{9000018107}
\Updated{2020-07-15 03:07:34}
\Level{B}
\SubArea{Linear equations and inequalities}
\LANGen{
\Question{
Solve the following inequality in the set of negative integers. 
\[ 
 \frac{x}
{6} + \frac{3x - 2} 
 {2} > -5
\]}
{\(x\in \left \{-2;-1\right \}\)}{\(x\in \left \{-3;-2;-1\right \}\)}{\(x\in \left \{-3;-2\right \}\)}{\(x\in \left \{-1\right \}\)}{}{}}
\LANGpl{
\Question{
Rozwiąż następującą nierówność w zbiorze liczb całkowitych ujemnych.
\[ 
 \frac{x}
{6} + \frac{3x - 2} 
 {2} > -5
\]}
{\(x\in \left \{-2;-1\right \}\)}{\(x\in \left \{-3;-2;-1\right \}\)}{\(x\in \left \{-3;-2\right \}\)}{\(x\in \left \{-1\right \}\)}{}{}}
\LANGcs{
\Question{
Určete všechna záporná celá čísla, která jsou řešením nerovnice:
\[\frac{x}
{6} + \frac{3x-2} 
 {2} > -5\]}
{\(x\in\left \{-2;-1\right \}\)}{\(x\in\left \{-3;-2;-1\right \}\)}{\(x\in\left \{-3;-2\right \}\)}{\(x\in\left \{-1\right \}\)}{}{}}
\LANGsk{
\Question{
Nájdite všetky záporné celé čísla, ktoré sú riešením nerovnice.
\[ 
 \frac{x}
{6} + \frac{3x - 2} 
 {2} > -5
\]}
{\(x\in \left \{-2;-1\right \}\)}{\(x\in \left \{-3;-2;-1\right \}\)}{\(x\in \left \{-3;-2\right \}\)}{\(x\in \left \{-1\right \}\)}{}{}}
\LANGes{
\Question{
Determina todos los enteros negativos que son solución de la siguiente inecuación. 
\[ 
 \frac{x}
{6} + \frac{3x - 2} 
 {2} > -5
\]}
{\(x\in \left \{-2;-1\right \}\)}{\(x\in \left \{-3;-2;-1\right \}\)}{\(x\in \left \{-3;-2\right \}\)}{\(x\in \left \{-1\right \}\)}{}{}}
\End

\Data\ID{9000021701}
\Updated{2020-07-15 03:07:34}
\Level{B}
\SubArea{Linear equations and inequalities}
\LANGen{
\Question{
Assume \(x\in [ - 2;2] \) 
and solve the following inequality. 
\[ 
 10 + 7x\leq 5 - 3x
\]}
{\(x\in \left [ -2;-\frac{1}
{2}\right ] \)}{\(x\in \left (-\infty ;-\frac{1}
{2}\right ] \)}{\(x\in \left [ -\frac{1}
{2};2\right ] \)}{\(x\in [ - 2;2] \)}{}{}}
\LANGpl{
\Question{
Zakładając, że  \(x\in [ - 2;2] \) 
rozwiąż podaną nierówność.
\[ 
 10 + 7x\leq 5 - 3x
\]}
{\(x\in \left [ -2;-\frac{1}
{2}\right ] \)}{\(x\in \left (-\infty ;-\frac{1}
{2}\right ] \)}{\(x\in \left [ -\frac{1}
{2};2\right ] \)}{\(x\in [ - 2;2] \)}{}{}}
\LANGcs{
\Question{
Vyberte všechna řešení nerovnice v
intervalu \(\langle - 2;2\rangle \).
\[10 + 7x\leq 5 - 3x\]}
{\(x\in\left \langle -2;-\frac{1}
{2}\right \rangle \)}{\(x\in\left (-\infty ;-\frac{1}
{2}\right \rangle \)}{\(x\in\left \langle -\frac{1}
{2};2\right \rangle \)}{\(x\in\langle - 2;2\rangle \)}{}{}}
\LANGsk{
\Question{
Vyberte všetky riešení nasledujúce nerovnice v intervale \(x\in [ - 2;2] \).
\[ 
 10 + 7x\leq 5 - 3x
\] je}
{\(x\in \left [ -2;-\frac{1}
{2}\right ] \)}{\(x\in \left (-\infty ;-\frac{1}
{2}\right ] \)}{\(x\in \left [ -\frac{1}
{2};2\right ] \)}{\(x\in [ - 2;2] \)}{}{}}
\LANGes{
\Question{
Suponiendo que \(x\in [ - 2;2] \) 
resuelve la siguiente inecuación: 
\[ 
 10 + 7x\leq 5 - 3x
\]}
{\(x\in \left [ -2;-\frac{1}
{2}\right ] \)}{\(x\in \left (-\infty ;-\frac{1}
{2}\right ] \)}{\(x\in \left [ -\frac{1}
{2};2\right ] \)}{\(x\in [ - 2;2] \)}{}{}}
\End

\Data\ID{9000021702}
\Updated{2020-07-15 03:07:34}
\Level{B}
\SubArea{Linear equations and inequalities}
\LANGen{
\Question{
Find the positive integer solutions of the following inequality. 
\[ 
 \frac{1 + x}
 {3} -\frac{8 - 3x} 
 {2} < \frac{3x} 
 {2} - 2
\]}
{\(x\in \{1;2;3;4\}\)}{\(x\in \mathbb{N}\)}{\(x\in \{1;2;3;4;5\}\)}{\(x\in [ 1;5] \)}{}{}}
\LANGpl{
\Question{
Znajdź naturalne liczbę rozwiązań podanej nierówności.
\[ 
 \frac{1 + x}
 {3} -\frac{8 - 3x} 
 {2} < \frac{3x} 
 {2} - 2
\]}
{\(x\in \{1;2;3;4\}\)}{\(x\in \mathbb{N}\)}{\(x\in \{1;2;3;4;5\}\)}{\(x\in [ 1;5] \)}{}{}}
\LANGcs{
\Question{
Vyřešte následující nerovnici v množině $\mathbb{N}$.
\[\frac{1+x}
 {3} -\frac{8-3x} 
 {2} < \frac{3x} 
 {2} - 2\]}
{\(x\in\{1;2;3;4\}\)}{\(x\in\mathbb{N}\)}{\(x\in\{1;2;3;4;5\}\)}{\(x\in\langle 1;5\rangle \)}{}{}}
\LANGsk{
\Question{
Vyriešte danú nerovnicu v množine $\mathbb{N}$.
\[ 
 \frac{1 + x}
 {3} -\frac{8 - 3x} 
 {2} < \frac{3x} 
 {2} - 2
\]}
{\(x\in \{1;2;3;4\}\)}{\(x\in \mathbb{N}\)}{\(x\in \{1;2;3;4;5\}\)}{\(x\in [ 1;5] \)}{}{}}
\LANGes{
\Question{
Resuelve la siguiente inecuación en el conjunto $\mathbb{N}$.
\[\frac{1+x}
 {3} -\frac{8-3x} 
 {2} < \frac{3x} 
 {2} - 2\]}
{\(x\in\{1;2;3;4\}\)}{\(x\in\mathbb{N}\)}{\(x\in\{1;2;3;4;5\}\)}{\(x\in\langle 1;5\rangle \)}{}{}}
\End

\Data\ID{9000021704}
\Updated{2020-07-15 03:07:34}
\Level{B}
\SubArea{Linear equations and inequalities}
\LANGen{
\Question{
Solve the following inequality. 
\[ 
 \frac{x + 1}
 {4} -\frac{x + 2} 
 {3} > \frac{x + 3} 
 {6} -\frac{3x - 4} 
 {12}
\]}
{\(x\in \emptyset \)}{\(x\in \mathbb{R}\)}{\(x\in (-\infty ;29)\)}{\(x\in \{0\}\)}{}{}}
\LANGpl{
\Question{
Rozwiąż podaną nierówność.
\[ 
 \frac{x + 1}
 {4} -\frac{x + 2} 
 {3} > \frac{x + 3} 
 {6} -\frac{3x - 4} 
 {12}
\]}
{\(x\in \emptyset \)}{\(x\in \mathbb{R}\)}{\(x\in (-\infty ;29)\)}{\(x\in \{0\}\)}{}{}}
\LANGcs{
\Question{
Vyřešte danou nerovnici.
\[\frac{x+1}
 {4} -\frac{x+2} 
 {3} > \frac{x+3} 
 {6} -\frac{3x-4} 
 {12} \]}
{\(x\in\emptyset \)}{\(x\in\mathbb{R}\)}{\(x\in(-\infty ;29)\)}{\(x\in\{0\}\)}{}{}}
\LANGsk{
\Question{
Vyriešte danú nerovnicu.
\[ 
 \frac{x + 1}
 {4} -\frac{x + 2} 
 {3} > \frac{x + 3} 
 {6} -\frac{3x - 4} 
 {12}
\]}
{\(x\in \emptyset \)}{\(x\in \mathbb{R}\)}{\(x\in (-\infty ;29)\)}{\(x\in \{0\}\)}{}{}}
\LANGes{
\Question{
Resuelve la siguiente inecuación. 
\[ 
 \frac{x + 1}
 {4} -\frac{x + 2} 
 {3} > \frac{x + 3} 
 {6} -\frac{3x - 4} 
 {12}
\]}
{\(x\in \emptyset \)}{\(x\in \mathbb{R}\)}{\(x\in (-\infty ;29)\)}{\(x\in \{0\}\)}{}{}}
\End

\Data\ID{9000021705}
\Updated{2020-07-15 03:07:34}
\Level{B}
\SubArea{Linear equations and inequalities}
\LANGen{
\Question{
Solve the following inequality in the set of negative integers. 
\[ 
 \frac{3x - 4}
 {2} -\frac{2x - 5} 
 {3} + \frac{3 - 4x} 
 {5} > 0
\]}
{\(x\in \{ - 7;-6;-5;-4;-3;-2;-1\}\)}{\(x\in \emptyset \)}{\(x\in [ - 8;0] \)}{\(x\in \{ - 8;-7;-6;-5;-4;-3;-2;-1\}\)}{}{}}
\LANGpl{
\Question{
Rozwiąż podaną nierówność w zbiorze liczb całkowitych ujemnych. 
\[ 
 \frac{3x - 4}
 {2} -\frac{2x - 5} 
 {3} + \frac{3 - 4x} 
 {5} > 0
\]}
{\(x\in \{ - 7;-6;-5;-4;-3;-2;-1\}\)}{\(x\in \emptyset \)}{\(x\in [ - 8;0] \)}{\(x\in \{ - 8;-7;-6;-5;-4;-3;-2;-1\}\)}{}{}}
\LANGcs{
\Question{
Vyřešte danou nerovnici v množině celých záporných čísel.
\[\frac{3x-4}
 {2} -\frac{2x-5} 
 {3} + \frac{3-4x} 
 {5} > 0\]}
{\(x\in\{ - 7;-6;-5;-4;-3;-2;-1\}\)}{\(x\in\emptyset \)}{\(x\in\left \langle -8;0\right \rangle \)}{\(x\in\{ - 8;-7;-6;-5;-4;-3;-2;-1\}\)}{}{}}
\LANGsk{
\Question{
Vyriešte danú nerovnicu v množine celých záporných čísel.
\[ 
 \frac{3x - 4}
 {2} -\frac{2x - 5} 
 {3} + \frac{3 - 4x} 
 {5} > 0
\]}
{\(x\in \{ - 7;-6;-5;-4;-3;-2;-1\}\)}{\(x\in \emptyset \)}{\(x\in [ - 8;0] \)}{\(x\in \{ - 8;-7;-6;-5;-4;-3;-2;-1\}\)}{}{}}
\LANGes{
\Question{
Resuelve la siguiente inecuación en el conjunto de los enteros negativos.
\[ 
 \frac{3x - 4}
 {2} -\frac{2x - 5} 
 {3} + \frac{3 - 4x} 
 {5} > 0
\]}
{\(x\in \{ - 7;-6;-5;-4;-3;-2;-1\}\)}{\(x\in \emptyset \)}{\(x\in [ - 8;0] \)}{\(x\in \{ - 8;-7;-6;-5;-4;-3;-2;-1\}\)}{}{}}
\End

\Data\ID{9000021709}
\Updated{2020-07-15 03:07:34}
\Level{B}
\SubArea{Linear equations and inequalities}
\LANGen{
\Question{
Find all the values of \(x\) 
for which the expression \(\frac{x+5}
 {4} -\frac{7-3x} 
 {12} \) 
is not bigger than \(\frac{2x+4}
 {6} + \frac{x-3} 
 {3} \).}
{\(x\in [ 6;\infty )\)}{\(x\in (6;\infty )\)}{\(x\in (-\infty ;6)\)}{\(x\in (-\infty ;6] \)}{}{}}
\LANGpl{
\Question{
Znajdź wszystkie wartości  \(x\), dla których wyrażenie \(\frac{x+5}
 {4} -\frac{7-3x} 
 {12} \) 
nie jest większe niż \(\frac{2x+4}
 {6} + \frac{x-3} 
 {3} \).}
{\(x\in [ 6;\infty )\)}{\(x\in (6;\infty )\)}{\(x\in (-\infty ;6)\)}{\(x\in (-\infty ;6] \)}{}{}}
\LANGcs{
\Question{
Výraz \(\frac{x+5}
 {4} -\frac{7-3x} 
 {12} \) nemá větší
hodnotu než výraz \(\frac{2x+4}
 {6} + \frac{x-3} 
 {3} \) 
pro:}
{\(x\in \left \langle 6;\infty \right )\)}{\(x\in (6;\infty )\)}{\(x\in (-\infty ;6)\)}{\(x\in \left (-\infty ;6\right \rangle \)}{}{}}
\LANGsk{
\Question{
Nájdite všetky hodnoty \(x\), pre ktoré výraz
 \(\frac{x+5}
 {4} -\frac{7-3x} 
 {12} \) 
nie je väčší ako výraz \(\frac{2x+4}
 {6} + \frac{x-3} 
 {3} \).}
{\(x\in [ 6;\infty )\)}{\(x\in (6;\infty )\)}{\(x\in (-\infty ;6)\)}{\(x\in (-\infty ;6] \)}{}{}}
\LANGes{
\Question{
La expresión \(\frac{x+5}
 {4} -\frac{7-3x} 
 {12} \) es menor que la expresión \(\frac{2x+4}
 {6} + \frac{x-3} 
 {3} \) 
para:}
{\(x\in \left \langle 6;\infty \right )\)}{\(x\in (6;\infty )\)}{\(x\in (-\infty ;6)\)}{\(x\in \left (-\infty ;6\right \rangle \)}{}{}}
\End

\Data\ID{9000021710}
\Updated{2020-07-15 03:07:34}
\Level{B}
\SubArea{Linear equations and inequalities}
\LANGen{
\Question{
Find the maximal integer \(x\) 
which satisfies the following inequality. 
\[ 
 \frac{x + 6} 
 {3} -\frac{x - 1} 
 {2} < 2 - 0.2x
\]}
{\(- 16\)}{\(- 15\)}{\(- 14\)}{\(14\)}{}{}}
\LANGpl{
\Question{
Znajdź największą liczbę całkowitą \(x\) 
która spełnia następującą nierówność.
\[ 
 \frac{x + 6} 
 {3} -\frac{x - 1} 
 {2} < 2 - 0.2x
\]}
{\(- 16\)}{\(- 15\)}{\(- 14\)}{\(14\)}{}{}}
\LANGcs{
\Question{
Určete největší celé číslo, které je řešením nerovnice: \[\frac{x+6}{3} -\frac{x-1}{2} < 2 - 0{,}2x\]}
{\(- 16\)}{\(- 15\)}{\(- 14\)}{\(14\)}{}{}}
\LANGsk{
\Question{
Určte najväčšie celé číslo, ktoré vyhovuje nerovnici:
\[\frac{x+6}
 {3} -\frac{x-1} 
 {2} < 2 - 0{,}2x\]}
{\(- 16\)}{\(- 15\)}{\(- 14\)}{\(14\)}{}{}}
\LANGes{
\Question{
Determina el número mayor \(x\) 
que es solución de la siguiente inecuación. 
\[ 
 \frac{x + 6} 
 {3} -\frac{x - 1} 
 {2} < 2 - 0.2x
\]}
{\(- 16\)}{\(- 15\)}{\(- 14\)}{\(14\)}{}{}}
\End

\Data\ID{9000021801}
\Updated{2020-07-15 03:07:34}
\Level{B}
\SubArea{Linear equations and inequalities}
\LANGen{
\Question{
Solve the following system of inequalities.
\[\begin{aligned}
 \frac{1}
{3}(2x + 5) &\geq 0.5\left (\frac{2 + 3x}
 {2} + 2\right ) & &
 \\0.2(3 - 2x) &\leq \frac{1}
{3}\left (\frac{4 - 2x}
 {5} + 2\right ) & &
 \end{aligned}\]}
{\(x\in \left [ -\frac{5}
{4};2\right ] \)}{\(x\in [ 2;\infty )\)}{\(x\in \left (-\infty ;-\frac{5}
{4}\right ] \)}{\(x\in \emptyset \)}{}{}}
\LANGpl{
\Question{
Rozwiąż następujący układ nierówności. 
\[\begin{aligned}
 \frac{1}
{3}(2x + 5) &\geq 0.5\left (\frac{2 + 3x}
 {2} + 2\right ) & &
 \\0.2(3 - 2x) &\leq \frac{1}
{3}\left (\frac{4 - 2x}
 {5} + 2\right ) & &
 \end{aligned}\]}
{\(x\in \left\langle -\frac{5}
{4};2\right \rangle \)}{\(x\in \langle 2;\infty )\)}{\(x\in \left (-\infty ;-\frac{5}
{4}\right \rangle \)}{\(x\in \emptyset \)}{}{}}
\LANGcs{
\Question{
Vyřešte následující soustavu nerovnic.
\[\begin{aligned}
 \frac{1}
{3}(2x + 5) &\geq 0{,}5\left (\frac{2 + 3x}
 {2} + 2\right ) & &
 \\0{,}2(3 - 2x) &\leq \frac{1}
{3}\left (\frac{4 - 2x}
 {5} + 2\right ) & &
 \end{aligned}\]}
{\(x\in\left \langle -\frac{5}
{4};2\right \rangle \)}{\(x\in\langle 2;\infty )\)}{\(x\in\left (-\infty ;-\frac{5}
{4}\right \rangle \)}{\(x\in\emptyset \)}{}{}}
\LANGsk{
\Question{
Vyriešte nasledujúcu sústavu nerovníc.
\[\begin{aligned}
 \frac{1}
{3}(2x + 5) &\geq 0.5\left (\frac{2 + 3x}
 {2} + 2\right ) & &
 \\0.2(3 - 2x) &\leq \frac{1}
{3}\left (\frac{4 - 2x}
 {5} + 2\right ) & &
 \end{aligned}\]}
{\(x\in \left \langle -\frac{5}
{4};2\right \rangle \)}{\(x\in \langle 2;\infty )\)}{\(x\in \left (-\infty ;-\frac{5}
{4}\right \rangle \)}{\(x\in \emptyset \)}{}{}}
\LANGes{
\Question{
Resuelve el siguiente sistema de inecuaciones.
\[\begin{aligned}
 \frac{1}
{3}(2x + 5) &\geq 0.5\left (\frac{2 + 3x}
 {2} + 2\right ) & &
 \\0.2(3 - 2x) &\leq \frac{1}
{3}\left (\frac{4 - 2x}
 {5} + 2\right ) & &
 \end{aligned}\]}
{\(x\in \left [ -\frac{5}
{4};2\right ] \)}{\(x\in [ 2;\infty )\)}{\(x\in \left (-\infty ;-\frac{5}
{4}\right ] \)}{\(x\in \emptyset \)}{}{}}
\End

\Data\ID{9000021802}
\Updated{2020-07-15 03:07:34}
\Level{B}
\SubArea{Linear equations and inequalities}
\LANGen{
\Question{
Solve the following system of inequalities.
\[\begin{aligned}
 15x - 2 &\geq 3x + 2 > 2x + 1 & &
 \\10x + 1 & > 5x + 1\geq 6 - x & &
 \end{aligned}\]}
{\(x\in \left [ \frac{5}
{6};\infty \right )\)}{\(x\in [ - 1;\infty )\)}{\(x\in \emptyset \)}{\(x\in [ 2;\infty )\)}{}{}}
\LANGpl{
\Question{
Rozwiąż następujący układ nierówności. 
\[\begin{aligned}
 15x - 2 &\geq 3x + 2 > 2x + 1 & &
 \\10x + 1 & > 5x + 1\geq 6 - x & &
 \end{aligned}\]}
{\(x\in \left \langle \frac{5}
{6};\infty \right )\)}{\(x\in \langle - 1;\infty )\)}{\(x\in \emptyset \)}{\(x\in \langle 2;\infty )\)}{}{}}
\LANGcs{
\Question{
Vyřešte následující soustavu nerovnic.
\[\begin{aligned}
 15x - 2 &\geq 3x + 2 > 2x + 1 & &
 \\10x + 1 & > 5x + 1\geq 6 - x & &
 \end{aligned}\]}
{\(x\in\left \langle \frac{5}
{6};\infty \right )\)}{\(x\in\langle - 1;\infty )\)}{\(x\in\emptyset \)}{\(x\in\langle 2;\infty )\)}{}{}}
\LANGsk{
\Question{
Vyriešte danú sústavu nerovníc.
\[\begin{aligned}
 15x - 2 &\geq 3x + 2 > 2x + 1 & &
 \\10x + 1 & > 5x + 1\geq 6 - x & &
 \end{aligned}\]}
{\(x\in \left \langle \frac{5}
{6};\infty \right )\)}{\(x\in \langle - 1;\infty )\)}{\(x\in \emptyset \)}{\(x\in \langle 2;\infty )\)}{}{}}
\LANGes{
\Question{
Resuelve el siguiente sistema de inecuaciones.
\[\begin{aligned}
 15x - 2 &\geq 3x + 2 > 2x + 1 & &
 \\10x + 1 & > 5x + 1\geq 6 - x & &
 \end{aligned}\]}
{\(x\in \left [ \frac{5}
{6};\infty \right )\)}{\(x\in [ - 1;\infty )\)}{\(x\in \emptyset \)}{\(x\in [ 2;\infty )\)}{}{}}
\End

\Data\ID{9000039002}
\Updated{2020-07-15 03:07:34}
\Level{B}
\SubArea{Linear equations and inequalities}
\LANGen{
\Question{
Assuming \(x\in \mathbb{Z}\),
find the solution set of the following system. 
\[ 
 -3\leq 2(x + 2)\leq 6
\]}
{\(\{ - 3;-2;-1;0;1\}\)}{\(\{ - 4;-2;-1;0;1\}\)}{\(\{ - 3;-2;-1;0\}\)}{\(\{0;1\}\)}{}{}}
\LANGpl{
\Question{
Zakładając, że \(x\in \mathbb{Z}\),
znajdź zbiór rozwiązań podanego układu.
\[ 
 -3\leq 2(x + 2)\leq 6
\]}
{\(\{ - 3;-2;-1;0;1\}\)}{\(\{ - 4;-2;-1;0;1\}\)}{\(\{ - 3;-2;-1;0\}\)}{\(\{0;1\}\)}{}{}}
\LANGcs{
\Question{
Je dána soustava nerovnic.
\[ 
 -3\leq 2(x + 2)\leq 6
\]
Množinou řešení této soustavy v
\(\mathbb{Z}\) je:}
{\(\{ - 3;-2;-1;0;1\}\)}{\(\{ - 4;-2;-1;0;1\}\)}{\(\{ - 3;-2;-1;0\}\)}{\(\{0;1\}\)}{}{}}
\LANGsk{
\Question{
Nájdite množinu riešenia danej sústavy nerovníc pre \(x\in \mathbb{Z}\). 
\[ 
 -3\leq 2(x + 2)\leq 6
\]}
{\(\{ - 3;-2;-1;0;1\}\)}{\(\{ - 4;-2;-1;0;1\}\)}{\(\{ - 3;-2;-1;0\}\)}{\(\{0;1\}\)}{}{}}
\LANGes{
\Question{
Suponiendo que \(x\in \mathbb{Z}\),
resuelve el siguiente sistema de inecuaciones.
\[ 
 -3\leq 2(x + 2)\leq 6
\]}
{\(\{ - 3;-2;-1;0;1\}\)}{\(\{ - 4;-2;-1;0;1\}\)}{\(\{ - 3;-2;-1;0\}\)}{\(\{0;1\}\)}{}{}}
\End

\Data\ID{9000039003}
\Updated{2020-07-15 03:07:34}
\Level{B}
\SubArea{Linear equations and inequalities}
\LANGen{
\Question{
Assuming \(x\in \mathbb{R}\),
find the solution set of the following system. 
\[ 
 x + 2 > 2x + 3 > 3x + 5
\]}
{\((-\infty ;-2)\)}{\((-2;+\infty )\)}{\((-\infty ;-1)\)}{\((-2;-1)\)}{}{}}
\LANGpl{
\Question{
Zakładając, że \(x\in \mathbb{R}\),
znajdź zbiór rozwiązań podanego układu.
\[ 
 x + 2 > 2x + 3 > 3x + 5
\]}
{\((-\infty ;-2)\)}{\((-2;+\infty )\)}{\((-\infty ;-1)\)}{\((-2;-1)\)}{}{}}
\LANGcs{
\Question{
Je dána soustava nerovnic.
\[ 
 x + 2 > 2x + 3 > 3x + 5
\]
Množinou řešení této soustavy v
\(\mathbb{R}\) je:}
{\((-\infty ;-2)\)}{\((-2;+\infty )\)}{\((-\infty ;-1)\)}{\((-2;-1)\)}{}{}}
\LANGsk{
\Question{
Nájdite množinu riešenia danej sústavy pre \(x\in \mathbb{R}\). 
\[ 
 x + 2 > 2x + 3 > 3x + 5
\]}
{\((-\infty ;-2)\)}{\((-2;+\infty )\)}{\((-\infty ;-1)\)}{\((-2;-1)\)}{}{}}
\LANGes{
\Question{
Suponiendo que \(x\in \mathbb{R}\),
resuelve el siguiente sistema de inecuaciones. 
\[ 
 x + 2 > 2x + 3 > 3x + 5
\]}
{\((-\infty ;-2)\)}{\((-2;+\infty )\)}{\((-\infty ;-1)\)}{\((-2;-1)\)}{}{}}
\End

\Data\ID{9100018007}
\Updated{2020-07-15 03:07:10}
\Level{B}
\SubArea{Linear equations and inequalities}
\LANGen{
\Question{
Graph the solution of the following inequality. 
\[ 
 -x < 3
\]}
{}{}{}{}{}{}}
\LANGpl{
\Question{
Przedstaw graficzne rozwiązanie podanej nierówności.
\[ 
 -x < 3
\]}
{}{}{}{}{}{}}
\LANGcs{
\Question{
Vyřešte nerovnici a vyberte správné řešení znázorněné na číselné ose. \[- x < 3\]}
{}{}{}{}{}{}}
\LANGsk{
\Question{
Vyriešte danú nerovnicu a vyberte správne riešenie znázornené na číselnej osi. 
\[ 
 -x < 3
\]}
{}{}{}{}{}{}}
\LANGes{
\Question{
NA}
{}{}{}{}{}{}}

\IMGanswer{000180_07-figure0.pdf}
\IMGanswer{000180_07-figure1.pdf}
\IMGanswer{000180_07-figure2.pdf}
\IMGanswer{000180_07-figure3.pdf}\End

\Data\ID{9100018008}
\Updated{2020-07-15 03:07:10}
\Level{B}
\SubArea{Linear equations and inequalities}
\LANGen{
\Question{
Graph the solution of the following inequality. 
\[ 
 \frac{x}
{2}\geq 3
\]}
{}{}{}{}{}{}}
\LANGpl{
\Question{
Który z poniższych rysunków przedstawia rozwiązanie następującej nierówności? 
\[ 
 \frac{x}
{2}\geq 3
\]}
{}{}{}{}{}{}}
\LANGcs{
\Question{
Vyřešte nerovnici a vyberte správné řešení znázorněné na číselné ose. \[\frac{x}
{2} \geq 3\]}
{}{}{}{}{}{}}
\LANGsk{
\Question{
Vyriešte danú nerovnicu a vyberte správne riešenie znázornené na číselnej osi. 
\[ 
 \frac{x}
{2}\geq 3
\]}
{}{}{}{}{}{}}
\LANGes{
\Question{
NA}
{}{}{}{}{}{}}

\IMGanswer{000180_08-figure1.pdf}
\IMGanswer{000180_08-figure0.pdf}
\IMGanswer{000180_08-figure2.pdf}
\IMGanswer{000180_08-figure3.pdf}\End

\Data\ID{9100018009}
\Updated{2020-07-15 03:07:10}
\Level{B}
\SubArea{Linear equations and inequalities}
\LANGen{
\Question{
Graph the solution of the following inequality. 
\[ 
 0\geq 3x
\]}
{}{}{}{}{}{}}
\LANGpl{
\Question{
Który z poniższych rysunków przedstawia rozwiązanie następującej nierówności? 
\[ 
 0\geq 3x
\]}
{}{}{}{}{}{}}
\LANGcs{
\Question{
Vyřešte danou nerovnici a vyberte správné řešení znázorněné na číselné ose. \[0\geq 3x\]}
{}{}{}{}{}{}}
\LANGsk{
\Question{
Riešte danú nerovnicu a vyberte správne riešenie znázornené na číselnej osi. 
\[ 
 0\geq 3x
\]}
{}{}{}{}{}{}}
\LANGes{
\Question{
NA}
{}{}{}{}{}{}}

\IMGanswer{000180_09-figure2.pdf}
\IMGanswer{000180_09-figure0.pdf}
\IMGanswer{000180_09-figure1.pdf}
\IMGanswer{000180_09-figure3.pdf}\End

\Data\ID{9100018108}
\Updated{2020-07-15 03:07:10}
\Level{B}
\SubArea{Linear equations and inequalities}
\LANGen{
\Question{
Graph the solution of the following inequality. 
\[ 
 \frac{x - 5}
 {2} \leq 2\left (x + 1\right )
\]}
{}{}{}{}{}{}}
\LANGpl{
\Question{
Który z poniższych rysunków przedstawia rozwiązanie następującej nierówności? 
\[ 
 \frac{x - 5}
 {2} \leq 2\left (x + 1\right )
\]}
{}{}{}{}{}{}}
\LANGcs{
\Question{
Vyřešte danou nerovnici a vyberte správné řešení znázorněné na číselné ose. \[\frac{x-5}
 {2} \leq 2\left (x + 1\right )\]}
{}{}{}{}{}{}}
\LANGsk{
\Question{
Riešte danú nerovnicu a vyberte správne riešenie znázornené na číselnej osi. 
\[ 
 \frac{x - 5}
 {2} \leq 2\left (x + 1\right )
\]}
{}{}{}{}{}{}}
\LANGes{
\Question{
NA}
{}{}{}{}{}{}}

\IMGanswer{000181_08-figure1.pdf}
\IMGanswer{000181_08-figure0.pdf}
\IMGanswer{000181_08-figure2.pdf}
\IMGanswer{000181_08-figure3.pdf}\End

\Data\ID{9100018109}
\Updated{2020-07-15 03:07:10}
\Level{B}
\SubArea{Linear equations and inequalities}
\LANGen{
\Question{
Graph the solution of the following inequality. 
\[ 
 1 -\frac{x - 4} 
 {5} \geq \frac{4x}
 {5}
\]}
{}{}{}{}{}{}}
\LANGpl{
\Question{
Który z poniższych rysunków przedstawia rozwiązanie następującej nierówności? 
\[ 
 1 -\frac{x - 4} 
 {5} \geq \frac{4x}
 {5}
\]}
{}{}{}{}{}{}}
\LANGcs{
\Question{
Vyřešte danou nerovnici a vyberte správné řešení znázorněné na číselné ose. \[1 -\frac{x-4} 
 {5} \geq \frac{4x}
 {5} \]}
{}{}{}{}{}{}}
\LANGsk{
\Question{
Riešte danú nerovnicu a vyberte správne riešenie znázornené na číselnej osi. 
\[ 
 1 -\frac{x - 4} 
 {5} \geq \frac{4x}
 {5}
\]}
{}{}{}{}{}{}}
\LANGes{
\Question{
NA}
{}{}{}{}{}{}}

\IMGanswer{000181_09-figure2.pdf}
\IMGanswer{000181_09-figure0.pdf}
\IMGanswer{000181_09-figure1.pdf}
\IMGanswer{000181_09-figure3.pdf}\End

\Data\ID{9100021805}
\Updated{2020-07-15 04:07:14}
\Level{B}
\SubArea{Linear equations and inequalities}
\LANGen{
\Question{
Identify a picture which shows the graphical solution of the following inequality in red.
\[ 
 \frac{1}
{2}x + 1\leq 3x - 4
\]}
{}{}{}{}{}{}}
\LANGpl{
\Question{
Który z poniższych rysunków przedstawia rozwiązanie następującej nierówności? Rozwiązanie jest wyświetlane na czerwono.
\[ 
 \frac{1}
{2}x + 1\leq 3x - 4
\]}
{}{}{}{}{}{}}
\LANGcs{
\Question{
Urči, na kterém obrázku je červenou barvou znázorněna množina všech řešení
nerovnice.
\[\frac{1}
{2}x + 1\leq 3x - 4\]}
{}{}{}{}{}{}}
\LANGsk{
\Question{
Urči, na ktorom obrázku je červenou farbou znázornená množina všetkých riešení danej nerovnice.
\[ 
 \frac{1}
{2}x + 1\leq 3x - 4
\]}
{}{}{}{}{}{}}
\LANGes{
\Question{
NA}
{}{}{}{}{}{}}

\IMGanswer{000218_05-figure0.pdf}
\IMGanswer{000218_05-figure1.pdf}
\IMGanswer{000218_05-figure2.pdf}
\IMGanswer{000218_05-figure3.pdf}\End

\Data\ID{1003029701}
\Updated{2020-07-15 03:07:12}
\Level{C}
\SubArea{Linear equations and inequalities}
\LANGen{
\Question{
One and a half of a loaf of bread weighs as same as a quarter of a loaf of bread together with a kilo weight. If it is possible, find how much one loaf weighs. (Assume that all the loaves weigh the same.)}
{\( 0.8\,\mathrm{kg} \)}{\( 1.25\,\mathrm{kg} \)}{\( 0.2\,\mathrm{kg} \)}{It is not possible to determine the exact weight of the loaf from the given information.}{}{}}
\LANGpl{
\Question{
Półtora bochenka chleba waży tyle samo, co jedna czwarta bochenka chleba razem z kilogramem. Jeśli to możliwe, oblicz ile waży jeden bochen. (Załóżmy, że wszystkie bochenki ważą tyle samo.)}
{\( 0{,}8\,\mathrm{kg} \)}{\( 1{,}25\,\mathrm{kg} \)}{\( 0{,}2\,\mathrm{kg} \)}{Podane informacje nie są wystarczające, by udzielić odpowiedzi.}{}{}}
\LANGcs{
\Question{
Jeden a půl bochníku chleba váží stejně jako čtvrt bochníku společně s kilogramovým závažím. Jestli to je možné, určete, kolik váží jeden bochník. (Předpokládáme, že všechny bochníky váží stejně).}
{\( 0{,}8\,\mathrm{kg} \)}{\( 1{,}25\,\mathrm{kg} \)}{\( 0{,}2\,\mathrm{kg} \)}{Z daných informací není možné jednoznačně určit hmotnost bochníku.}{}{}}
\LANGsk{
\Question{
Jeden a pol bochníka chleba váži rovnako ako štvrť bochníka spoločne s kilogramovým závažím. Ak je to možné, určte, koľko váži jeden bochník. (Predpokladáme, že všetky bochníky vážia rovnako.)}
{\( 0{,}8\,\mathrm{kg} \)}{\( 1{,}25\,\mathrm{kg} \)}{\( 0{,}2\,\mathrm{kg} \)}{Z daných informácií nie je možné určiť presnú hmotnosť bochníka.}{}{}}
\LANGes{
\Question{
Una barra y media de pan pesa igual que un cuarto de barra junto con un kilo de peso. Si es posible, determina cuánto pesa una barra de pan. (Se supone que todas las barras de pan pesan igual).}
{\( 0.8\,\mathrm{kg} \)}{\( 1.25\,\mathrm{kg} \)}{\( 0.2\,\mathrm{kg} \)}{A partir de la información dada, no es posible determinar el peso exacto de una barra de pan.}{}{}}
\End

\Data\ID{1003029702}
\Updated{2020-07-15 03:07:12}
\Level{C}
\SubArea{Linear equations and inequalities}
\LANGen{
\Question{
The comparative math test consists of the total of three problems.  From students taking the test \( 16\% \) solved correctly all three problems, \( 3/5 \) of the students made mistake in exactly one of the problems and \( 2/15 \) of the students made mistakes in exactly two of the problems. Eight students gave neither one of the problems. How many students wrote the test without a mistake?}
{\( 12 \)}{\( 75 \)}{\( 45 \)}{\( 10 \)}{}{}}
\LANGpl{
\Question{
Porównywalny test matematyczny składa się z trzech zagadnień. Wśród uczniów przystępujących do testu \( 16\% \)  poprawnie rozwiązało wszystkie trzy problemy, \( 3/5 \) uczniów popełniło błąd w jednym z zagadnień, a \( 2/15 \) uczniów popełniło błędy w dokładnie dwóch zagadnieniach. Ośmiu studentów nie udzieliło żadnej odpowiedzi. Ilu uczniów napisało test bez błędu?}
{\( 12 \)}{\( 75 \)}{\( 45 \)}{\( 10 \)}{}{}}
\LANGcs{
\Question{
Z výsledků srovnávacích testů z matematiky vyplývá, že \( 16\,\% \) studentů vypočítalo správně všechny tři úlohy. Právě jednu úlohu mělo špatně \( 3/5 \) studentů a \( 2/15 \) studentů mělo špatně právě dvě úlohy. Osm studentů nevyřešilo správně ani jednu úlohu. Kolik studentů napsalo bezchybný test?}
{\( 12 \)}{\( 75 \)}{\( 45 \)}{\( 10 \)}{}{}}
\LANGsk{
\Question{
Porovnávací test z matematiky pozostával z troch úloh. Všetky tri úlohy vypočítalo správne \( 16\% \) študentov, \( 3/5 \) študentov vypočítalo nesprávne práve jednu úlohu a \( 2/15 \) študentov vypočítalo nesprávne práve dve úlohy. Osem študentov nevyriešilo správne ani jednu úlohu. Koľko študentov napísalo test bez chyby?}
{\( 12 \)}{\( 75 \)}{\( 45 \)}{\( 10 \)}{}{}}
\LANGes{
\Question{
La prueba de nivel de matemáticas consta, en total, de tres ejercicios. De todos los estudiantes, el 16% resolvieron correctamente los tres ejercicios, 3/5 de los estudiantes cometieron errores exactamente en uno de los ejercicios y 2/15 de los estudiantes cometieron errores en dos de todos los ejercicios. Ocho estudiantes no resolvieron ninguno de los ejercicios. ¿Cuántos estudiantes escribieron la prueba sin cometer errores?}
{\( 12 \)}{\( 75 \)}{\( 45 \)}{\( 10 \)}{}{}}
\End

\Data\ID{1003029703}
\Updated{2020-07-15 03:07:12}
\Level{C}
\SubArea{Linear equations and inequalities}
\LANGen{
\Question{
A sample of \( 168 \) respondents of a sociological research splits into three categories according to nationality: Czech, Slovak and other. There is four times less Slovak respondents than Czech in the sample. The other nationality reported \( 85\% \) less respondents than Czech nationality. Choose the equation that represents the algebraic relation between the nationalities of the sample and leads to evaluation of the number of people of the given nationality.}
{\( x+4x+\frac{4x\cdot15}{100}=168 \), where \( x \) is the number of people of the Slovak nationality}{\( x+\frac x4+\frac{x\cdot15}{4\cdot100}=168 \), where \( x \) is the number of people of the Slovak nationality}{\( x+4x+\frac{4x\cdot85}{100}=168 \), where \( x \) is the number of people of the Slovak nationality}{\( x+4x+\frac{4x\cdot15}{100}=168 \), where \( x \) is the number of people of the Czech nationality}{}{}}
\LANGpl{
\Question{
Grupa \( 168 \)  respondentów badania socjologicznego dzieli się na trzy kategorie według narodowości: czeska, słowacka i inna. W grupie jest cztery razy mniej respondentów słowackich niż czeskich. Inna narodowość zgłosiła \( 85\% \)  mniej respondentów niż narodowość czeska. Wybierz równanie, które przedstawia algebraiczną zależność między narodowościami tej grupy i pozwoli określić  liczbę osób danej narodowości.}
{\( x+4x+\frac{4x\cdot15}{100}=168 \), gdzie \( x \) stanowi liczbę osób narodowości słowackiej.}{\( x+\frac x4+\frac{x\cdot15}{4\cdot100}=168 \), gdzie \( x \) stanowi liczbę osób narodowości słowackiej.}{\( x+4x+\frac{4x\cdot85}{100}=168 \), gdzie \( x \) stanowi liczbę osób narodowości słowackiej.}{\( x+4x+\frac{4x\cdot15}{100}=168 \), gdzie \( x \) stanowi liczbę osób narodowości czeskiej.}{}{}}
\LANGcs{
\Question{
Sociologického průzkumu se zúčastnilo celkem \( 168 \) respondentů. Z jeho výsledků plyne, že část respondentů je české národnosti, čtyřikrát méně lidí je slovenské národnosti a k některé další národnosti se přihlásilo o \( 85\,\% \) méně lidí než k české. Vyberte rovnici, která představuje algebraické řešení dané slovní úlohy a vede ke zjištění počtu lidí dané národnosti.}
{\( x+4x+\frac{4x\cdot15}{100}=168 \), kde \( x \) je počet lidí slovenské národnosti}{\( x+\frac x4+\frac{x\cdot15}{4\cdot100}=168 \), kde \( x \) je počet lidí slovenské národnosti}{\( x+4x+\frac{4x\cdot85}{100}=168 \), kde \( x \) je počet lidí slovenské národnosti}{\( x+4x+\frac{4x\cdot15}{100}=168 \), kde \( x \) je počet lidí české národnosti}{}{}}
\LANGsk{
\Question{
Sociologického prieskumu sa zúčastnilo celkom \( 168 \) respondentov. Z jeho výsledkov vyplýva, že časť respondentov je českej národnosti, slovenskej národnosti je štyrikrát menej ako českej a k niektorej ďalšej národnosti sa prihlásilo o \( 85\% \) menej respondentov než k českej. Vyberte rovnicu, ktorá predstavuje algebraické riešenie danej slovnej úlohy a vedie k zisteniu počtu ľudí danej národnosti.}
{\( x+4x+\frac{4x\cdot15}{100}=168 \), kde \( x \) je počet ľudí slovenskej národnosti}{\( x+\frac x4+\frac{x\cdot15}{4\cdot100}=168 \), kde \( x \) je počet ľudí slovenskej národnosti}{\( x+4x+\frac{4x\cdot85}{100}=168 \), kde \( x \) je počet ľudí slovenskej národnosti}{\( x+4x+\frac{4x\cdot15}{100}=168 \), kde \( x \) je počet ľudí českej národnosti}{}{}}
\LANGes{
\Question{
Una muestra de \( 168 \) encuestados de una investigación sociológica se divide en tres categorías según la nacionalidad: checa, eslovaca y otras. Hay cuatro veces menos encuestados eslovacos que checos en la muestra. Otras nacionalidades forman un \( 85\% \) menos que los de nacionalidad checa. Elige la ecuación que representa la relación algebraica entre las nacionalidades de la muestra y determina el número de personas de las nacionalidades dadas.}
{\( x+4x+\frac{4x\cdot15}{100}=168 \), donde \( x \) es el número de personas de nacionalidad eslovaca}{\( x+\frac x4+\frac{x\cdot15}{4\cdot100}=168 \), donde \( x \) es el número de personas de nacionalidad eslovaca}{\( x+4x+\frac{4x\cdot85}{100}=168 \), donde \( x \) es el número de personas de nacionalidad eslovaca}{\( x+4x+\frac{4x\cdot15}{100}=168 \), donde \( x \) es el número de personas de nacionalidad checa}{}{}}
\End

\Data\ID{1003029704}
\Updated{2020-07-15 03:07:12}
\Level{C}
\SubArea{Linear equations and inequalities}
\LANGen{
\Question{
Tom is \( 15 \) years old and his grandpa is \( 67 \). How many years will it take for the grandpa to be three times as old as Tom? Choose the equation which represents the algebraic solution to the given word problem.}
{\( 67+x=3\cdot(15+x) \)}{\( 67=3\cdot x \)}{\( 67=3\cdot(15+x) \)}{\( 3\cdot(67+x)=3\cdot(15+x) \)}{}{}}
\LANGpl{
\Question{
Tomek ma  \( 15 \) lat, a jego dziadek \( 67 \). Za ile lat dziadek będzie trzy razy starszy od Tomka? Wybierz równanie, które przedstawia algebraiczne rozwiązanie tego problemu.}
{\( 67+x=3\cdot(15+x) \)}{\( 67=3\cdot x \)}{\( 67=3\cdot(15+x) \)}{\( 3\cdot(67+x)=3\cdot(15+x) \)}{}{}}
\LANGcs{
\Question{
Tomášovi je \( 15 \) roků a jeho dědovi je \( 67 \). Za kolik roků může Tomáš tvrdit, že děda je třikrát starší než on? Vyberte rovnici, která představuje algebraické řešení dané slovní úlohy.}
{\( 67+x=3\cdot(15+x) \)}{\( 67=3\cdot x \)}{\( 67=3\cdot(15+x) \)}{\( 3\cdot(67+x)=3\cdot(15+x) \)}{}{}}
\LANGsk{
\Question{
Tomáš má \( 15 \) rokov a jeho starý otec má \( 67 \) rokov. O koľko rokov bude starý otec trikrát starší než Tomáš? Vyberte rovnicu, ktorá predstavuje algebraické riešenie danej slovnej úlohy.}
{\( 67+x=3\cdot(15+x) \)}{\( 67=3\cdot x \)}{\( 67=3\cdot(15+x) \)}{\( 3\cdot(67+x)=3\cdot(15+x) \)}{}{}}
\LANGes{
\Question{
Tomás tiene  \( 15 \) años y su abuelo tiene  \( 67 \). ¿Dentro de cuántos años la edad del abuelo será el triple que la de Tomás? Elige la ecuación que representa la relación algebraica del ejercicio puesto.}
{\( 67+x=3\cdot(15+x) \)}{\( 67=3\cdot x \)}{\( 67=3\cdot(15+x) \)}{\( 3\cdot(67+x)=3\cdot(15+x) \)}{}{}}
\End

\Data\ID{1003029705}
\Updated{2020-07-15 03:07:12}
\Level{C}
\SubArea{Linear equations and inequalities}
\LANGen{
\Question{
In a tank, there are \( 60 \) litres of \( 2\% \) saline solution. How much water must evaporate to increase salt concentration into \( 3\% \)?}
{\( 20 \) litres}{\( 40 \) litres}{\( 1.8 \) litres}{\( 58.2 \) litres}{}{}}
\LANGpl{
\Question{
W zbiorniku znajduje się \( 60 \) litrów \( 2\% \) roztworu soli fizjologicznej. Ile wody należy odparować, żeby zwiększyć stężenie soli do \( 3\% \)?}
{\( 20 \) litrów}{\( 40 \) litrów}{\( 1{,}8 \) litrów}{\( 58{,}2 \) litrów}{}{}}
\LANGcs{
\Question{
V nádobě je \( 60 \) litrů \( 2\,\% \) roztoku kuchyňské soli. Kolik vody se musí vypařit, aby se koncentrace zvětšila na \( 3\,\% \)?}
{\( 20 \) litrů}{\( 40 \) litrů}{\( 1{,}8 \) litrů}{\( 58{,}2 \) litrů}{}{}}
\LANGsk{
\Question{
V nádobe je \( 60 \) litrov  \( 2\% \) roztoku kuchynskej soli. Koľko vody sa musí vypariť, aby sa koncentrácia soli zvýšila na \( 3\% \)?}
{\( 20 \) litrov}{\( 40 \) litrov}{\( 1{,}8 \) litra}{\( 58{,}2 \) litra}{}{}}
\LANGes{
\Question{
Un tanque contiene  \( 60 \) litros de solución salina al  \( 2\% \) ¿Qué cantidad de agua debe evaporarse para aumentar la concentración de sal al  \( 3\% \)?}
{\( 20 \) litros}{\( 40 \) litros}{\( 1.8 \) litros}{\( 58.2 \) litros}{}{}}
\End

\Data\ID{1003029706}
\Updated{2020-07-15 03:07:12}
\Level{C}
\SubArea{Linear equations and inequalities}
\LANGen{
\Question{
In a river, water flows at speed of \( 1\,\mathrm{mps} \). A boat, which moves at speed of \( 4\,\mathrm{mps} \) in a calm water, carries mail to a little town distant \( 6\,\mathrm{km} \) downstream. How long will it take till the boat returns? (We disregard the time necessary for mail-handover.) Note: Abbreviation "mps" stands for meter per second.}
{\( 53\,\mathrm{min}\ 20\,\mathrm{s} \)}{\( 50\,\mathrm{min} \)}{\( 3\,\mathrm{min}\ 12\,\mathrm{s} \)}{\( 1\,\mathrm{min}\ 20\,\mathrm{s} \)}{}{}}
\LANGpl{
\Question{
W rzece woda płynie z prędkością \( 1\,\mathrm{mps} \). Łódź, która porusza się z prędkością 4 m / s w spokojnej wodzie, przenosi pocztę do małego miasteczka odległego o \( 6\,\mathrm{km} \)  w dół rzeki.  Ile czasu minie, zanim łódź wróci? (Nie uwzględniamy czasu potrzebnego na przekazanie poczty.) Uwaga: Skrót "mps" oznacza metr na sekundę.}
{\( 53\,\mathrm{min}\ 20\,\mathrm{s} \)}{\( 50\,\mathrm{min} \)}{\( 3\,\mathrm{min}\ 12\,\mathrm{s} \)}{\( 1\,\mathrm{min}\ 20\,\mathrm{s} \)}{}{}}
\LANGcs{
\Question{
Voda v řece plyne rychlostí \( 1\,\mathrm{m/s} \). Člun, který se na klidné vodě pohybuje rychlostí \( 4\,\mathrm{m/s} \), vezl poštu do městečka vzdáleného \( 6\,\mathrm{km} \) po proudu. Jak dlouho potrvá, než se vrátí zpět? (Dobu potřebnou na předání pošty zanedbáváme.)}
{\( 53\,\mathrm{min}\ 20\,\mathrm{s} \)}{\( 50\,\mathrm{min} \)}{\( 3\,\mathrm{min}\ 12\,\mathrm{s} \)}{\( 1\,\mathrm{min}\ 20\,\mathrm{s} \)}{}{}}
\LANGsk{
\Question{
Voda v rieke tečie rýchlosťou \( 1\,\mathrm{m/s} \). Čln, ktorý sa na pokojnej vode pohybuje rýchlosťou \( 4\,\mathrm{m/s} \), vezie poštu do mestečka vzdialeného \( 6\,\mathrm{km} \) po prúde. Ako dlho potrvá, než sa čln vráti späť? (Dobu potrebnú na odovzdanie pošty zanedbáme.)}
{\( 53\,\mathrm{min}\ 20\,\mathrm{s} \)}{\( 50\,\mathrm{min} \)}{\( 3\,\mathrm{min}\ 12\,\mathrm{s} \)}{\( 1\,\mathrm{min}\ 20\,\mathrm{s} \)}{}{}}
\LANGes{
\Question{
En un río el agua fluye a una velocidad de  \( 1\,\mathrm{mps} \). Un bote, que se mueve con velocidad de \( 4\,\mathrm{mps} \) en aguas tranquilas, lleva el correo y se desplaza hasta un pueblo pequeño situado a \( 6\,\mathrm{km} \) río abajo. ¿Cuánto tiempo emplea el bote en su recorrido de vuelta? (No contamos el tiempo necesario para la entrega del correo).}
{\( 53\,\mathrm{min}\ 20\,\mathrm{s} \)}{\( 50\,\mathrm{min} \)}{\( 3\,\mathrm{min}\ 12\,\mathrm{s} \)}{\( 1\,\mathrm{min}\ 20\,\mathrm{s} \)}{}{}}
\End

\Data\ID{1003031101}
\Updated{2020-07-15 03:07:12}
\Level{C}
\SubArea{Linear equations and inequalities}
\LANGen{
\Question{
So far Johnny has received the grades: \( 5 \), \( 3 \), \( 3 \), \( 3 \), \( 2 \), \( 2 \), \( 1 \), \( 1 \), for the math course in this school term. What another grade he needs to get so that the average of grades is lower than \( 2.5 \)?
(Assumption: All grades are of the same weight in grading scale of \( 1 \), \( 2 \), \( 3 \), \( 4 \), \( 5 \), where \( 1 \) is the best result.)}
{at worst \( 2 \)}{at worst \( 3 \)}{only \( 1 \)}{The arithmetic mean can’t be less than \( 2.5 \) in any case.}{}{}}
\LANGpl{
\Question{
Do tej pory Janek otrzymał oceny: \( 5 \), \( 3 \), \( 3 \), \( 3 \), \( 2 \), \( 2 \), \( 1 \), \( 1 \),  za kursy matematyczne w tym półroczu. Jaką kolejną ocenę musi uzyskać, aby średnia ocen była niższa niż \( 2{,}5 \)  ? (Założenie: wszystkie oceny mają tę samą wagę w skali ocen \( 1 \), \( 2 \), \( 3 \), \( 4 \), \( 5 \), gdzie \( 1 \) jest najlepszym wynikiem.)}
{w najgorszym przypadku \( 2 \)}{w najgorszym przypadku \( 3 \)}{tylko  \( 1 \)}{Średnia arytmetyczna nie może być niższa niż \( 2{,}5 \) w żadnym przypadku.}{}{}}
\LANGcs{
\Question{
Honza zatím dostal v tomto pololetí tyto známky z matematiky:  \( 5 \), \( 3 \), \( 3 \), \( 3 \), \( 2 \), \( 2 \), \( 1 \), \( 1 \). Jakou musí dostat poslední známku, aby aritmetický průměr za pololetí byl lepší než 2,5? (Předpokládáme, že všechny známky mají stejnou váhu a platí pětistupňová klasifikační stupnice.)}
{nejhůře \( 2 \)}{nejhůře \( 3 \)}{pouze \( 1 \)}{Průměr nebude v žádném případě lepší než \( 2{,}5 \).}{}{}}
\LANGsk{
\Question{
Ján zatiaľ dostal v tomto polroku tieto známky z matematiky: \( 5 \), \( 3 \), \( 3 \), \( 3 \), \( 2 \), \( 2 \), \( 1 \), \( 1 \). Akú ďalšiu známku musí dostať, aby aritmetický priemer za polrok bol lepší než \( 2{,}5 \)?
(Predpokladáme, že všetky známky majú rovnakú váhu a platí päťstupňová klasifikačná stupnica: \( 1 \), \( 2 \), \( 3 \), \( 4 \), \( 5 \), kde \( 1 \) je najlepšia známka.)}
{najhoršie \( 2 \)}{najhoršie \( 3 \)}{len \( 1 \)}{Aritmetický priemer nebude v žiadnom prípade lepší než \( 2{,}5 \).}{}{}}
\LANGes{
\Question{
Hasta ahora Juan ha sacado en matemáticas este trimestre las notas siguientes: \( 5 \), \( 3 \), \( 3 \), \( 3 \), \( 2 \), \( 2 \), \( 1 \), \( 1 \). ¿Qué otra nota necesita sacar para que el promedio de calificaciones sea inferior a \( 2.5 \)?
(Se supone que todas las notas tienen el mismo peso en la escala de calificación de  \( 1 \), \( 2 \), \( 3 \), \( 4 \), \( 5 \), siendo \( 1 \) la mejor nota).}
{a lo peor \( 2 \)}{a lo peor \( 3 \)}{sólo \( 1 \)}{La media aritmética no puede ser inferior a  \( 2.5 \) en ningún caso.}{}{}}
\End

\Data\ID{1003031103}
\Updated{2020-07-15 03:07:12}
\Level{C}
\SubArea{Linear equations and inequalities}
\LANGen{
\Question{
Five litres of a quality wine in by yourself provided bottles cost more than three and a half litres of the same wine in a bigger wine container. The container price is \( 150\,\mathrm{CZK} \). Complete the next statement so that it is true. The price of one litre of this quality wine is}
{greater than \( 100\,\mathrm{CZK} \).}{less than \( 100\,\mathrm{CZK} \).}{greater than \( 350\,\mathrm{CZK} \).}{greater than \( 500\,\mathrm{CZK} \).}{}{}}
\LANGpl{
\Question{
Pięć litrów wina gatunkowego w butelkach zapewnianych na własną rękę  kosztuje ponad trzy i pół litra tego samego wina w większym pojemniku na wino. Cena pojemnika wynosi \( 150\,\mathrm{CZK} \). Dokończ podane stwierdzenie, aby było prawdziwe. Cena jednego litra tego wina gatunkowego jest}
{wyższa niż \( 100\,\mathrm{CZK} \).}{niższa niż  \( 100\,\mathrm{CZK} \).}{wyższa niż  \( 350\,\mathrm{CZK} \).}{wyższa niż \( 500\,\mathrm{CZK} \).}{}{}}
\LANGcs{
\Question{
\( 5 \) litrů kvalitního vína ve vlastních nádobách stojí více než \( 3{,}5 \) litrů téhož vína v demižónu, jehož cena je  \( 150\,\mathrm{CZK} \). Dokončete následující tvrzení tak, aby bylo pravdivé. Cena jednoho litru tohoto vína je}
{vyšší než \( 100\,\mathrm{CZK} \).}{nižší než \( 100\,\mathrm{CZK} \).}{vyšší než \( 350\,\mathrm{CZK} \).}{vyšší než \( 500\,\mathrm{CZK} \).}{}{}}
\LANGsk{
\Question{
Päť litrov kvalitného vína vo vlastných nádobách stojí viac než tri a pol litra tohto vína v demižóne, ktorého cena \( 150\,\mathrm{CZK} \). Dokončite nasledujúce tvrdenie tak, aby bolo pravdivé. Cena jedného litra tohto vína je}
{vyššia než \( 100\,\mathrm{CZK} \).}{nižšia než \( 100\,\mathrm{CZK} \).}{vyššia než \( 350\,\mathrm{CZK} \).}{vyššia než \( 500\,\mathrm{CZK} \).}{}{}}
\LANGes{
\Question{
Cinco litros de vino de calidad vendidos en botellas de plástico cuestan más que tres litros y medio del mismo vino vendidos en una garrafa que cuesta 150 coronas checas. Completa la siguiente expresión para que sea verdadera. El precio de un litro de este vino de calidad es}
{más de \( 100\,\mathrm{CZK} \).}{menos de \( 100\,\mathrm{CZK} \).}{más de \( 350\,\mathrm{CZK} \).}{más de \( 500\,\mathrm{CZK} \).}{}{}}
\End

\Data\ID{1003031104}
\Updated{2020-07-15 03:07:12}
\Level{C}
\SubArea{Linear equations and inequalities}
\LANGen{
\Question{
Dan and Jane took a bike trip. Dan rode for \( 3 \) hours at constant speed. Jane rode for half an hour longer at the speed of \( 4\,\mathrm{kph} \) less than Dan’s speed. Identify, which of the following statements about Dan’s speed is true.}
{The speed is less than \( 28\,\mathrm{kph} \).}{The speed is greater than \( 28\,\mathrm{kph} \).}{The speed is less than \( 20\,\mathrm{kph} \).}{The speed is greater than \( 24\,\mathrm{kph} \).}{}{}}
\LANGpl{
\Question{
Daniel i Janek wybrali się na wycieczkę rowerową. Daniel jechał przez \(3 \) godzin ze stałą prędkością. Janek jechał o pół godziny dłużej z prędkością \(4 \, \mathrm {kph} \) mniejszą  niż prędkość Daniela. Określ, które z poniższych stwierdzeń na temat prędkości Daniela jest prawdziwe?}
{Prędkość jest niższa niż  \( 28\,\mathrm{kph} \).}{Prędkość jest wyższa niż \( 28\,\mathrm{kph} \).}{Prędkość jest niższa niż \( 20\,\mathrm{kph} \).}{Prędkość jest wyższa niż \( 24\,\mathrm{kph} \).}{}{}}
\LANGcs{
\Question{
Na cyklistickém výletě jel Dan  \( 3 \) hodiny stálou rychlostí a urazil při tom větší vzdálenost než Jana, která sice jela o půl hodiny déle, ale její rychlost byla o  \( 4\,\mathrm{km/h} \) nižší než Danova.  Určete, který z následujících výroků o Danově rychlosti je pravdivý.}
{Rychlost je nižší než \( 28\,\mathrm{km/h} \).}{Rychlost je vyšší než \( 28\,\mathrm{km/h} \).}{Rychlost je vyšší než \( 20\,\mathrm{km/h} \).}{Rychlost je vyšší než \( 24\,\mathrm{km/h} \).}{}{}}
\LANGsk{
\Question{
Daniel a Jana išli na cyklistický výlet. Daniel išiel \( 3 \) hodiny stálou rýchlosťou. Jana išla o pol hodinu dlhšie, ale jej rýchlosť bola o \( 4\,\mathrm{km/h} \) nižšia než Danielova rýchlosť. Určte, ktorý z nasledujúcich výrokov o Danielovej rýchlosti je pravdivý.}
{Rýchlosť je nižšia než \( 28\,\mathrm{km/h} \).}{Rýchlosť je vyššia než \( 28\,\mathrm{km/h} \).}{Rýchlosť je nižšia než \( 20\,\mathrm{km/h} \).}{Rýchlosť je vyššia než \( 24\,\mathrm{km/h} \).}{}{}}
\LANGes{
\Question{
Daniel y Juana viajaron en bicicleta. Daniel pedaleó \( 3 \) horas a una velocidad constante. Juana pedaleó media hora más que Daniel pero a una velocidad de \( 4\,\mathrm{kph} \) inferior. Identifica cuál de las siguientes afirmaciones sobre la velocidad de Daniel es verdadera.}
{La velocidad es menor que \( 28\,\mathrm{kph} \).}{La velocidad es mayor que \( 28\,\mathrm{kph} \).}{La velocidad es menor que \( 20\,\mathrm{kph} \).}{La velocidad es mayor que \( 24\,\mathrm{kph} \).}{}{}}
\End

\Data\ID{1003197401}
\Updated{2020-07-15 03:07:21}
\Level{C}
\SubArea{Linear equations and inequalities}
\LANGen{
\Question{
A man on a bike rides to a distant city at an average speed of \( 24\,\mathrm{kph} \). He will end the trip \( 12 \) minutes earlier, if he increases his average speed by \( 1\,\mathrm{kph} \). How distant is the city?}
{\( 120\,\mathrm{km} \)}{\( 115.2\,\mathrm{km} \)}{\( 300\,\mathrm{km} \)}{\( 125\,\mathrm{km} \)}{}{}}
\LANGpl{
\Question{
Mężczyzna jedzie na rowerze do odległego miasta ze średnią prędkością  \( 24\,\mathrm{km}/\mathrm{h} \). Zakończy podróż \( 12 \) minut wcześniej jeśli zwiększy swoją średnią prędkość o \( 1\,\mathrm{km}/\mathrm{h} \). Jak daleko znajduje się miasto?}
{\( 120\,\mathrm{km} \)}{\( 115{,}2\,\mathrm{km} \)}{\( 300\,\mathrm{km} \)}{\( 125\,\mathrm{km} \)}{}{}}
\LANGcs{
\Question{
Cyklista jede do vzdáleného města průměrnou rychlostí \( 24\,\mathrm{km}/\mathrm{h} \). Jestli zvýší svou průměrnou rychlost o \( 1 \,\mathrm{km}/\mathrm{h} \), dorazí do cíle o \( 12 \) minut dříve. Jak daleko je jeho cíl?}
{\( 120\,\mathrm{km} \)}{\( 115{,}2\,\mathrm{km} \)}{\( 300\,\mathrm{km} \)}{\( 125\,\mathrm{km} \)}{}{}}
\LANGsk{
\Question{
Cyklista ide do vzdialeného mesta priemernou rýchlosťou \( 24\,\mathrm{km}/\mathrm{h} \). Dorazí do cieľa o \( 12 \) minút skôr, ak zvýši svoju priemernú rýchlosť  o \( 1\,\mathrm{km}/\mathrm{h} \). Ako ďaleko je jeho cieľ?}
{\( 120\,\mathrm{km} \)}{\( 115{,}2\,\mathrm{km} \)}{\( 300\,\mathrm{km} \)}{\( 125\,\mathrm{km} \)}{}{}}
\LANGes{
\Question{
Un hombre monta en bicicleta para ir a una ciudad distinta a una velocidad media de \( 24\,\mathrm{kph} \). Si aumenta la velocidad media en \( 1\,\mathrm{kph} \), llegará a la ciudad \( 12 \) minutos antes. ¿A qué distancia está la de ciudad?}
{\( 120\,\mathrm{km} \)}{\( 115.2\,\mathrm{km} \)}{\( 300\,\mathrm{km} \)}{\( 125\,\mathrm{km} \)}{}{}}
\End

\Data\ID{1003197402}
\Updated{2020-07-15 03:07:21}
\Level{C}
\SubArea{Linear equations and inequalities}
\LANGen{
\Question{
Paul rides the bike at a constant speed of \( 18\,\mathrm{kph} \). Eighteen minutes after Paul starts his trip, Tom takes the same route on a motorbike at an average speed of \( 40\,\mathrm{kph} \). How far behind Paul will Tom be after \( 12 \) minutes of his ride?}
{\( 1\,\mathrm{km} \)}{\( 60\,\mathrm{km} \)}{\( 14\,\mathrm{km} \)}{after \( 12 \) minutes of ride Tom will be in front of Paul}{}{}}
\LANGpl{
\Question{
Paweł jedzie na rowerze ze stałą prędkością  \( 18\,\mathrm{km}/\mathrm{h} \). Osiemnaście minut po tym, jak Paweł rozpoczął swoją podróż, Tomek wyruszył tą samą trasą na motorze poruszając się ze średnią prędkością  \( 40\,\mathrm{km}/\mathrm{h} \). Jak daleko za Pawłem będzie Tomek po \( 12 \) minutach swojej podróży?}
{\( 1\,\mathrm{km} \)}{\( 60\,\mathrm{km} \)}{\( 14\,\mathrm{km} \)}{po \( 12 \) minutach jazdy  Tomek wyprzedzi Pawła}{}{}}
\LANGcs{
\Question{
Pavel jede na kole stálou rychlostí \( 18\,\mathrm{km}/\mathrm{h} \). Za \( 18 \) minut za ním po stejné trase vyjede Tomáš na motorce průměrnou rychlostí \( 40\,\mathrm{km}/\mathrm{h} \). Jak daleko za Pavlem bude Tomáš po \( 12 \) minutách jízdy?}
{\( 1\,\mathrm{km} \)}{\( 60\,\mathrm{km} \)}{\( 14\,\mathrm{km} \)}{Po \( 12 \) minutách jízdy bude Tomáš před Pavlem.}{}{}}
\LANGsk{
\Question{
Pavol jazdí na bicykli stálou rýchlosťou \( 18\,\mathrm{km}/\mathrm{h}\). O osemnásť  minút za ním po rovnakej trase vyrazí  Tomáš na motorke priemernou rýchlosťou \( 40\,\mathrm{km}/\mathrm{h} \). Ako ďaleko za Pavlom bude Tomáš po \( 12 \) minútach jazdy?}
{\( 1\,\mathrm{km} \)}{\( 60\,\mathrm{km} \)}{\( 14\,\mathrm{km} \)}{Po \( 12 \) minútach jazdy bude Tomáš pred Pavlom}{}{}}
\LANGes{
\Question{
Pablo monta en bicicleta a una velocidad constante de \( 18\,\mathrm{kph} \). Después de 18 minutos de que Pablo comenzara su viaje, Tom realiza la misma ruta en una motocicleta a una velocidad media de \( 40\,\mathrm{kph} \). ¿A cuántos kilómetros estará Tom detrás de Pablo después de \( 12 \) minutos de viaje?}
{\( 1\,\mathrm{km} \)}{\( 60\,\mathrm{km} \)}{\( 14\,\mathrm{km} \)}{Después de  \( 12 \) minutos de viaje, Tom estará frente a Pablo}{}{}}
\End

\Data\ID{1003197403}
\Updated{2020-07-15 03:07:21}
\Level{C}
\SubArea{Linear equations and inequalities}
\LANGen{
\Question{
An express train of length \( 150\,\mathrm{m} \) travels at the constant speed of \( 144\,\mathrm{kph} \). On a parallel track in the opposite direction travels a freight train of length \( 240\,\mathrm{m} \) at the constant speed of \( 90\,\mathrm{kph} \). How long does it take for trains to pass each other?}
{\( 6\,\mathrm{s} \)}{\( 1.\overline{6}\,\mathrm{s} \)}{\( 7.\overline{2} \)}{\( 26\,\mathrm{s} \)}{}{}}
\LANGpl{
\Question{
Pociąg pośpieszny o długości \( 150\,\mathrm{m} \)  porusza się ze stałą prędkością \( 144\,\mathrm{km}/\mathrm{h} \). Na równoległym torze w przeciwnym kierunku porusza się pociąg towarowy o długości \( 240\,\mathrm{m} \) z tą samą stałą prędkością \( 90\,\mathrm{km}/\mathrm{h} \). Ile czasu zajmie pociągom minięcie się?}
{\( 6\,\mathrm{s} \)}{\( 1{,}\overline{6}\,\mathrm{s} \)}{\( 7{,}\overline{2} \)}{\( 26\,\mathrm{s} \)}{}{}}
\LANGcs{
\Question{
Rychlík jede stálou rychlostí \( 144\,\mathrm{km}/\mathrm{h} \) a míjí se s protijedoucím nákladním vlakem, který jede rychlostí \( 90\,\mathrm{km}/\mathrm{h} \). Jak dlouho trvá míjení obou vlaků?  Víme, že rychlík je dlouhý \( 150\,\mathrm{m} \) a nákladní vlak je dlouhý \( 240\,\mathrm{m} \).}
{\( 6\,\mathrm{s} \)}{\( 1{,}\overline{6}\,\mathrm{s} \)}{\( 7{,}\overline{2} \)}{\( 26\,\mathrm{s} \)}{}{}}
\LANGsk{
\Question{
Rýchlik dlhý \( 150\,\mathrm{m} \) ide stálou rýchlosťou of \( 144\,\mathrm{km}/\mathrm{h} \) a stretáva sa s protiidúcim nákladným vlakom, ktorý je dlhý \( 240\,\mathrm{m} \) a ide stálou rýchlosťou \( 90\,\mathrm{km}/\mathrm{h} \). Ako dlho trvá míňanie sa vlakov?}
{\( 6\,\mathrm{s} \)}{\( 1{,}\overline{6}\,\mathrm{s} \)}{\( 7{,}\overline{2} \)}{\( 26\,\mathrm{s} \)}{}{}}
\LANGes{
\Question{
Un tren expreso, de \( 150\,\mathrm{m} \) metros de longitud, viaja a una velocidad constante de  \( 144\,\mathrm{kph} \). En una vía paralela en sentido opuesto viaja un tren de carga, de \( 240\,\mathrm{m} \) metros de longitud, a una velocidad constante de \( 90\,\mathrm{kph} \). ¿Cuánto tiempo tardarán en cruzarse?}
{\( 6\,\mathrm{s} \)}{\( 1.\overline{6}\,\mathrm{s} \)}{\( 7.\overline{2} \)}{\( 26\,\mathrm{s} \)}{}{}}
\End

\Data\ID{1003197404}
\Updated{2020-07-15 03:07:21}
\Level{C}
\SubArea{Linear equations and inequalities}
\LANGen{
\Question{
The stock of the monitored company lost \( 12\,\% \) of its value during a week. Its fall continued over the next week, and the value declined by another \( 4\,\% \). Let’s denote the original value of the stock by an \( x \). From the following possibilities, choose the expression that gives the value of the stock at the end of the monitored term.}
{\( 0.96\cdot0.88x \)}{\( (0.96+0.88)x \)}{\( 0.04\cdot0.12x \)}{\( [1-(0.04+0.12)]x \)}{}{}}
\LANGpl{
\Question{
Towar monitorowanej firmy stracił \( 12\,\% \) swojej wartości w ciągu tygodnia. Jego spadek trwał przez kolejny tydzień i wartość zmniejszyła się o kolejne \( 4\,\% \). Załóżmy, że początkowa wartość towaru oznaczona zostanie  za pomocą \( x \). Z podanych możliwości wybierz wyrażenie, które przedstawia wartość towaru na koniec monitorowania.}
{\( 0{,}96\cdot0{,}88x \)}{\( (0{,}96+0{,}88)x \)}{\( 0{,}04\cdot0{,}12x \)}{\( [1-(0{,}04+0{,}12)]x \)}{}{}}
\LANGcs{
\Question{
Akcie sledovaného podniku ztratily během týdne \( 12\,\% \) své hodnoty. Jejich pád ale dál pokračoval a během následujícího týdne se jejich hodnota snížila o další \( 4\,\% \). Označme \( x \) původní hodnotu akcií. Z nabídnutých možností vyberte výraz, pomocí kterého určíte hodnotu akcií na konci sledovaného období.}
{\( 0{,}96\cdot0{,}88x \)}{\( (0{,}96+0{,}88)x \)}{\( 0{,}04\cdot0{,}12x \)}{\( [1-(0{,}04+0{,}12)]x \)}{}{}}
\LANGsk{
\Question{
Akcie sledovaného podniku stratili počas týždňa \( 12\,\% \) svojej hodnoty. Ich pád ďalej pokračoval a počas nasledujúceho týždňa sa ich hodnota znížila o \( 4\,\% \). Označme \( x \) pôvodnú hodnotu akcií. Z ponúknutých možností vyberte výraz, pomocou ktorého určíte hodnotu akcií na konci sledovaného obdobia.}
{\( 0{,}96\cdot0{,}88x \)}{\( (0{,}96+0{,}88)x \)}{\( 0{,}04\cdot0{,}12x \)}{\( [1-(0{,}04+0{,}12)]x \)}{}{}}
\LANGes{
\Question{
Las acciones de una compañía perdieron \( 12\,\% \) de su valor durante una semana. Su caída continuó también durante la semana siguiente disminuyendo su valor en otro \( 4\,\% \). Denotemos el valor original de las acciones por una \( x \). De las siguientes posibilidades, determina la expresión que da el valor de las acciones al final del proceso.}
{\( 0.96\cdot0.88x \)}{\( (0.96+0.88)x \)}{\( 0.04\cdot0.12x \)}{\( [1-(0.04+0.12)]x \)}{}{}}
\End

\Data\ID{1003197405}
\Updated{2020-07-15 03:07:21}
\Level{C}
\SubArea{Linear equations and inequalities}
\LANGen{
\Question{
Nine people travel by bus. The same number of people gets off the bus at each of three bus stops and then so many people get in to double the number of remaining bus passengers. After the third bus stop there are \( 30 \) people in the bus. How many passengers get off at each bus stop?}
{\( 3 \)}{\( 2 \)}{\( 1 \)}{\( 6 \)}{}{}}
\LANGpl{
\Question{
Dziewięć osób podróżuje autobusem. Ta sama liczba osób wysiada z autobusu na każdym z trzech przystanków, a następnie wsiada tyle osób, że liczba pozostałych pasażerów autobusu podwaja się. Po trzecim przystanku w autobusie pozostało \( 30 \) osób w autobusie. Ilu pasażerów wysiada na każdym przystanku?}
{\( 3 \)}{\( 2 \)}{\( 1 \)}{\( 6 \)}{}{}}
\LANGcs{
\Question{
Autobusem cestuje \( 9 \) lidí. Na každé ze tří zastávek vystoupí stejný počet lidí a pak jich nastoupí tolik, aby se počet lidí, kteří v autobuse zůstanou po výstupu, zdvojnásobil. Po třetí zastávce jede v autobuse \( 30 \) lidí. Kolik pasažérů na každé zastávce vystupuje?}
{\( 3 \)}{\( 2 \)}{\( 1 \)}{\( 6 \)}{}{}}
\LANGsk{
\Question{
Autobusom cestuje deväť ľudí. Na každej z troch zastávok vystúpi rovnaký počet ľudí a potom ich nastúpi toľko, aby sa počet ľudí, ktorí v autobuse zostanú po výstupe, zdvojnásobil. Po tretej zastávke cestuje v autobuse \( 30 \) ľudí. Koľko pasažierov na každej zastávke vystupuje?}
{\( 3 \)}{\( 2 \)}{\( 1 \)}{\( 6 \)}{}{}}
\LANGes{
\Question{
Nueve personas viajan en un autobús. En cada una de las tres paradas baja la misma cantidad de personas y luego sube una cantidad para que se duplique el número de pasajeros que ha quedado en el autobús. Después de la tercera parada, en el autobús viajan \( 30 \) personas. ¿Cuántos pasajeros bajan en cada parada?}
{\( 3 \)}{\( 2 \)}{\( 1 \)}{\( 6 \)}{}{}}
\End

\Data\ID{1003197406}
\Updated{2020-07-15 03:07:21}
\Level{C}
\SubArea{Linear equations and inequalities}
\LANGen{
\Question{
Each of two companies should deliver the same amount of raw material. While inspecting it was found out that the company \( A \) delivered \( 150\,\mathrm{kg} \) and the company \( B \) delivered \( 194\,\mathrm{kg} \). At the moment of inspection, company \( A \) has to deliver yet three times more than what is left to be delivered by company \( B \). Choose the equation which does NOT correspond to the described situation.}
{\( 3(x-150)=x-194 \), where \( x \) represents the total planned delivery of both companies.}{\( x-150=3(x-194) \), where \( x \) represents the total planned delivery of both companies.}{\( 150+3x=194+x \), where \( x \) represents the amount of raw material which is left to be delivered by company \( B \).}{\( 150+x=194+\frac x3 \), where \( x \) represents the amount of raw material which is left to be delivered by company \( A \).}{}{}}
\LANGpl{
\Question{
Każda z dwóch firm powinna dostarczać taką samą ilość surowca. Podczas inspekcji wykryto, że firma \( A \) dostarczyła \( 150\,\mathrm{kg} \), a firma \( B \) dostarczyła \( 194\,\mathrm{kg} \). W chwili inspekcji firma \( A \) musi jeszcze dostarczyć trzy razy więcej surowca niż ilość, która pozostała do dostarczenia firmie \( B \). Wybierz równanie, które  NIE odnosi się do opisanej sytuacji.}
{\( 3(x-150)=x-194 \), gdzie \( x \) reprezentuje całkowitą zaplanowaną dostawę obydwu firm.}{\( x-150=3(x-194) \), gdzie \( x \) reprezentuje całkowitą zaplanowaną dostawę obydwu firm.}{\( 150+3x=194+x \), gdzie \( x \) reprezentuje ilość surowca, który pozostał do dostarczenia firmie \( B \).}{\( 150+x=194+\frac x3 \), gdzie \( x \) reprezentuje ilość surowca, który pozostał do dostarczenia firmie \( A \).}{}{}}
\LANGcs{
\Question{
Každá ze dvou firem má dodat stejné množství surovin. Při kontrole se zjistilo, že firma \( A \) dodala \( 150\,\mathrm{kg} \) a firma \( B \) dodala \( 194\,\mathrm{kg} \) suroviny. V okamžiku kontroly musí firma \( A \) dodat ještě trojnásobek toho, co zbývá dodat firmě \( B \). Z následujících rovnic vyberte takovou, která NENÍ matematickým vyjádřením popsaného stavu.}
{\( 3(x-150)=x-194 \), kde \( x \) vyjadřuje celkovou plánovanou dodávku obou podniků.}{\( x-150=3(x-194) \), kde \( x \) vyjadřuje celkovou plánovanou dodávku obou podniků.}{\( 150+3x=194+x \), kde \( x \) vyjadřuje množství surovin, které podnik \( B \) ještě musí dodat.}{\( 150+x=194+\frac x3 \), kde \( x \) vyjadřuje množství surovin, které podnik \( A \) ještě musí dodat.}{}{}}
\LANGsk{
\Question{
Každá z dvoch firiem má dodať rovnaké množstvo surovín. Pri kontrole sa zistilo, že  firma \( A \) dodala \( 150\,\mathrm{kg} \) a firma \( B \) dodala \( 194\,\mathrm{kg} \) suroviny. V čase kontroly musí firma \( A \) dodať ešte trojnásobok toho, čo zostáva dodať firme \( B \). Z nasledujúcich rovníc vyberte takú, ktorá NIE JE matematickým vyjadrením opísanej situácie.}
{\( 3(x-150)=x-194 \), kde \( x \) vyjadruje celkovú plánovanú dodávku obidvoch firiem}{\( x-150=3(x-194) \), kde \( x \) vyjadruje celkovú plánovanú dodávku obidvoch firiem}{\( 150+3x=194+x \), kde \( x \) vyjadruje množstvo surovín, ktoré firma \( B \) ešte musí dodať}{\( 150+x=194+\frac x3 \), kde \( x \) vyjadruje množstvo surovín, ktoré firma \( A \) ešte musí dodať}{}{}}
\LANGes{
\Question{
Dos empresas \( A \) y \( B \) deben suministrar la misma cantidad de materias primas. Al inspeccionarlas se descubrió que la empresa \( A \) suministró \( 150\,\mathrm{kg} \) y la empresa \( B \) suministró \( 194\,\mathrm{kg} \). En el momento de la inspección, la empresa \( A \) debe suministrar tres veces más de lo que queda por suministrar por la parte de la empresa \( B \). De las siguientes ecuaciones, determina cuál no corresponde a la situación descrita.}
{\( 3(x-150)=x-194 \), donde \( x \) representa el suministro total planificado de ambas empresas.}{\( x-150=3(x-194) \), donde \( x \) representa el suministro total planificado de ambas empresas.}{\( 150+3x=194+x \), donde \( x \) representa la cantidad de materias primas que le queda por suministrar a la empresa \( B \).}{\( 150+x=194+\frac x3 \), donde \( x \) representa la cantidad de materias primas que le queda por suministrar a la empresa  \( A \).}{}{}}
\End

\Data\ID{1003197407}
\Updated{2020-07-15 03:07:21}
\Level{C}
\SubArea{Linear equations and inequalities}
\LANGen{
\Question{
A hose can fill the garden pool in \( 20 \) hours. The owners bought a pump that guarantees filling the pool in \( 8 \) hours. However, a crack appeared in the pool's shell which can drain all the water from the full pool within \( 5 \) days. How long will it take to fill the pool with the hose and the pump at the same time as the water runs off through the crack?}
{exactly \( 6 \) hours}{approximately \( 5.7 \) hours}{approximately \( 5.5 \) hours}{approximately \( 6.8 \) hours}{}{}}
\LANGpl{
\Question{
Wąż może wypełnić basen ogrodowy w  \( 20 \) godzin. Właściciele kupili pompę, za pomocą której basen napełni się w \( 8 \) godzin. Jednakże w  muszli basenu pojawiło się pęknięcie, które sprawi, że wypełniony basen opróżni się w ciągu  \( 5 \) dni. Ile czasu zajmie napełnienie basenu za pomocą węża i pompy jeżeli woda wypływa przez pęknięcie?}
{dokładnie \( 6 \) godzin}{około \( 5{,}7 \) godziny}{około \( 5{,}5 \) godziny}{około \( 6{,}8 \) godziny}{}{}}
\LANGcs{
\Question{
Zahradní bazén se napouštěl hadicí \( 20 \) hodin. Majitelé dokoupili čerpadlo, kterým je garantováno napuštění bazénu za \( 8 \) hodin. V plášti bazénu ale vznikla trhlinka, kterou by mohla veškerá voda z plného bazénu odtéct během \( 5 \) dnů. Jak dlouho bude trvat napouštění hadicí i čerpadlem současně, když voda bude trhlinou odtékat?}
{Přesně \( 6 \) hodin.}{Přibližně \( 5{,}7 \) hodiny.}{Přibližně \( 5{,}5 \) hodiny.}{Přibližně \( 6{,}8 \) hodiny.}{}{}}
\LANGsk{
\Question{
Záhradný bazén sa napúšťa hadicou \( 20 \) hodín. Majitelia kúpili čerpadlo, ktorým je garantované napustenie bazénu za \( 8 \) hodín. Avšak, v plášti bazénu vznikla trhlina, cez ktorú by mohla všetka voda z plného bazénu odtiecť za \( 5 \) dní. Ako dlho bude trvať napustenie bazéna hadicou a čerpadlom súčasne, keď voda bude trhlinou odtekať?}
{presne \( 6 \) hodín}{približne \( 5{,}7 \) hodín}{približne \( 5{,}5 \) hodín}{približne \( 6{,}8 \) hodín}{}{}}
\LANGes{
\Question{
Una manguera es capaz de llenar una piscina de jardín en \( 20 \) horas. Los propietarios de la piscina compraron una bomba que garantiza que llena la piscina en \( 8 \) horas. Sin embargo, en el revestimiento de la piscina apareció una grieta que podría vaciar toda el agua de la piscina en \( 5 \) días. ¿Cuánto tiempo tardará el llenado de la piscina con la manguera y la bomba al mismo tiempo, si el agua va perdiendo por la grieta?}
{exactamente \( 6 \) horas}{aproximadamente \( 5.7 \) horas}{aproximadamente \( 5.5 \) horas}{aproximadamente \( 6.8 \) horas}{}{}}
\End

\Data\ID{1003197408}
\Updated{2020-07-15 03:07:21}
\Level{C}
\SubArea{Linear equations and inequalities}
\LANGen{
\Question{
The painter \( A \) takes \( 15 \) hours to paint a flat. The painter \( B \) takes \( 12 \) hours to paint the same flat and the painter \( C \) takes \( 10 \) hours to complete the same task. The painter \( A \) starts to paint alone and after \( 2 \) hours joins him the painter \( B \). After the next hour the painter \( C \) starts to paint with them too. How long will all three painters paint together before completing the work?}
{\( 2\,\mathrm{h}\ 52\,\mathrm{min} \)}{\( 3\,\mathrm{h}\ 8\,\mathrm{min} \)}{\( 3\,\mathrm{h}\ 24\,\mathrm{min} \)}{\( 4\,\mathrm{h} \)}{}{}}
\LANGpl{
\Question{
Malarzowi \( A \) malowanie mieszkania zajmuje \( 15 \) godzin. Malarzowi \( B \) pomalowanie tego samego mieszkania zajmuje \( 12 \) godzin, zaś malarzowi  \( C \)  \( 10 \) godzin. Malarz \( A \) zaczyna malować sam, a po  \( 2 \) godzinach dołącza do niego malarz \( B \). Po kolejnej godzinie malarz \( C \) zaczyna malować razem z nimi.  Jak długo malarze będą malować razem zanim ukończą zadanie?}
{\( 2\,\mathrm{h}\ 52\,\mathrm{min} \)}{\( 3\,\mathrm{h}\ 8\,\mathrm{min} \)}{\( 3\,\mathrm{h}\ 24\,\mathrm{min} \)}{\( 4\,\mathrm{h} \)}{}{}}
\LANGcs{
\Question{
Malíř \( A \) vymaluje celý byt za \( 15 \) hodin. Malíř \( B \) vymaluje stejný byt za \( 12 \) hodin a malíři \( C \) by stejný úkol trval \( 10 \) hodin. Daný byt začne malovat malíř \( A \), po \( 2 \) hodinách se k němu přidá malíř \( B \) a za další hodinu se přidá i malíř \( C \). Jak dlouho budou pracovat všichni společně, než dokončí práci?}
{\( 2\,\mathrm{h}\ 52\,\mathrm{min} \)}{\( 3\,\mathrm{h}\ 8\,\mathrm{min} \)}{\( 3\,\mathrm{h}\ 24\,\mathrm{min} \)}{\( 4\,\mathrm{h} \)}{}{}}
\LANGsk{
\Question{
Maliar \( A \) vymaľuje celý byt za \( 15 \) hodín. Maliar \( B \) vymaľuje ten istý byt za \( 12 \) hodín a maliarovi \( C \) by rovnaká úloha trvala \( 10 \) hodín. Daný byt začne maľovať maliar \( A \), po \( 2 \) hodinách sa k nemu pridá maliar \( B \) a za ďalšiu hodinu sa pridá aj maliar \( C \). Ako dlho budú pracovať všetci spoločne, než dokončia prácu?}
{\( 2\,\mathrm{h}\ 52\,\mathrm{min} \)}{\( 3\,\mathrm{h}\ 8\,\mathrm{min} \)}{\( 3\,\mathrm{h}\ 24\,\mathrm{min} \)}{\( 4\,\mathrm{h} \)}{}{}}
\LANGes{
\Question{
El pintor \( A \) tarda \( 15 \) horas en pintar un piso, el pintor \( B \) tarda \( 12 \) horas en pintar el mismo piso y el pintor \( C \) tarda \( 10 \) horas en completar la misma tarea. El pintor \( A \) empieza a pintar solo y después de \( 2 \) horas comienza a pintar también el pintor \( B \). Después de la siguiente hora, el pintor \( C \) empieza a pintar con ellos también. ¿Cuánto tiempo tardarán en terminar el trabajo desde que pintan todos juntos?}
{\( 2\,\mathrm{h}\ 52\,\mathrm{min} \)}{\( 3\,\mathrm{h}\ 8\,\mathrm{min} \)}{\( 3\,\mathrm{h}\ 24\,\mathrm{min} \)}{\( 4\,\mathrm{h} \)}{}{}}
\End

\Data\ID{9000018010}
\Updated{2020-07-15 03:07:34}
\Level{C}
\SubArea{Linear equations and inequalities}
\LANGen{
\Question{
The salary of Peter has been increased by
\(\$2\: 400\). The salary of
Jane increased by \(3\, \%\) 
and this increase has been bigger than the increase of the Peter's salary. In the
following list identify a possible old salary of Jane.}
{\(\$81\: 000\)}{\(\$80\: 000\)}{\(\$9\: 000\)}{\(\$8\: 000\)}{}{}}
\LANGpl{
\Question{
Pensja Piotra wzrosła o 
\(\$2\: 400\). Pensja Janki wzrosła o \(3\, \%\) 
i była to podwyżka większa niż podwyżka, którą otrzymał Piotr. Jaka była pensja Janki przed podwyżką?}
{\(\$81\: 000\)}{\(\$80\: 000\)}{\(\$9\: 000\)}{\(\$8\: 000\)}{}{}}
\LANGcs{
\Question{
Petrovi zvýšili plat o \(2\: 400\) Kč.
Lence zvýšili plat pouze o $3\,\%$ a přesto bylo její zvýšení větší než
Petrovo. Jaký může být Lenčin původní plat?}
{\(81\: 000\) Kč}{\(80\: 000\) Kč}{\(9\: 000\) Kč}{\(8\: 000\) Kč}{}{}}
\LANGsk{
\Question{
Petrovi zvýšili plat o 
\(\$2\: 400\). Lenke zvýšili plat o \(3\, \%\) 
a toto zvýšenie bolo väčšie než Petrovo. Aký môže byť Lenkin pôvodný plat?}
{\(\$81\: 000\)}{\(\$80\: 000\)}{\(\$9\: 000\)}{\(\$8\: 000\)}{}{}}
\LANGes{
\Question{
El salario de Pedro ha aumentado en \(\$2\: 400\). El salario de Juana ha aumentado en un  \(3\, \%\). Esa subida ha sido mayor que la subida del salario de Pedro. ¿Cuál era el salario original de Juana?}
{\(\$81\: 000\)}{\(\$80\: 000\)}{\(\$9\: 000\)}{\(\$8\: 000\)}{}{}}
\End

\Data\ID{9000018110}
\Updated{2020-07-15 03:07:34}
\Level{C}
\SubArea{Linear equations and inequalities}
\LANGen{
\Question{
Consider a coil composed from a ferrite core and a wire winding. The mass of the
core is \(2\, \mathrm{kg}\).
The material for the wire is such that the mass of the wire of the length
\(30\, \mathrm{m}\) is bigger than the mass
of the wire of the length \(10\, \mathrm{m}\) 
together with the ferrite core. In the following list identify a possible mass of one
meter of the wire.}
{\(110\, \mathrm{g}\)}{\(100\, \mathrm{g}\)}{\(0.01\, \mathrm{kg}\)}{\(0.09\, \mathrm{kg}\)}{}{}}
\LANGpl{
\Question{
Rozważmy cewkę składającą się z rdzenia ferrytowego i zwoju drutu. Masa rdzenia jest równa  \(2\, \mathrm{kg}\).
Materiał na drut jest taki, że masa drutu o długości 
\(30\, \mathrm{m}\) jest większa niż
masa drutu o długości \(10\, \mathrm{m}\) 
razem z ferrytowym rdzeniem. Jaka będzie masa 1 metra drutu?}
{\(110\, \mathrm{g}\)}{\(100\, \mathrm{g}\)}{\(0.01\, \mathrm{kg}\)}{\(0.09\, \mathrm{kg}\)}{}{}}
\LANGcs{
\Question{
Cívka na měděný drát má hmotnost \(2\, \mathrm{kg}\). \(30\, \mathrm{m}\) drátu bez cívky má větší hmotnost než \(10\, \mathrm{m}\) drátu s cívkou. Jaká může být hmotnost jednoho metru drátu?}
{\(110\, \mathrm{g}\)}{\(100\, \mathrm{g}\)}{\(0{,}01\, \mathrm{kg}\)}{\(0{,}09\, \mathrm{kg}\)}{}{}}
\LANGsk{
\Question{
Cievka na medený drôt má hmotnosť
\(2\, \mathrm{kg}\).
\(30\, \mathrm{m}\) drôtu bez cievky má
väčšiu hmotnosť ako \(10\, \mathrm{m}\) 
drôtu s cievkou. Aká  môže byť hmotnosť jedného metra drôtu?}
{\(110\, \mathrm{g}\)}{\(100\, \mathrm{g}\)}{\(0{,}01\, \mathrm{kg}\)}{\(0{,}09\, \mathrm{kg}\)}{}{}}
\LANGes{
\Question{
Un carrete para alambre de cobre pesa \(2\, \mathrm{kg}\). \(30\, \mathrm{m}\) de alambre sin carrete pesan más que \(10\, \mathrm{m}\) de alambre con carrete. ¿Cuánto pesa un metro de alambre?}
{\(110\, \mathrm{g}\)}{\(100\, \mathrm{g}\)}{\(0.01\, \mathrm{kg}\)}{\(0.09\, \mathrm{kg}\)}{}{}}
\End
